\documentclass[aspectratio=169]{beamer}\usepackage[]{graphicx}\usepackage[]{xcolor}
% maxwidth is the original width if it is less than linewidth
% otherwise use linewidth (to make sure the graphics do not exceed the margin)
\makeatletter
\def\maxwidth{ %
  \ifdim\Gin@nat@width>\linewidth
    \linewidth
  \else
    \Gin@nat@width
  \fi
}
\makeatother

\definecolor{fgcolor}{rgb}{0.345, 0.345, 0.345}
\newcommand{\hlnum}[1]{\textcolor[rgb]{0.686,0.059,0.569}{#1}}%
\newcommand{\hlsng}[1]{\textcolor[rgb]{0.192,0.494,0.8}{#1}}%
\newcommand{\hlcom}[1]{\textcolor[rgb]{0.678,0.584,0.686}{\textit{#1}}}%
\newcommand{\hlopt}[1]{\textcolor[rgb]{0,0,0}{#1}}%
\newcommand{\hldef}[1]{\textcolor[rgb]{0.345,0.345,0.345}{#1}}%
\newcommand{\hlkwa}[1]{\textcolor[rgb]{0.161,0.373,0.58}{\textbf{#1}}}%
\newcommand{\hlkwb}[1]{\textcolor[rgb]{0.69,0.353,0.396}{#1}}%
\newcommand{\hlkwc}[1]{\textcolor[rgb]{0.333,0.667,0.333}{#1}}%
\newcommand{\hlkwd}[1]{\textcolor[rgb]{0.737,0.353,0.396}{\textbf{#1}}}%
\let\hlipl\hlkwb

\usepackage{framed}
\makeatletter
\newenvironment{kframe}{%
 \def\at@end@of@kframe{}%
 \ifinner\ifhmode%
  \def\at@end@of@kframe{\end{minipage}}%
  \begin{minipage}{\columnwidth}%
 \fi\fi%
 \def\FrameCommand##1{\hskip\@totalleftmargin \hskip-\fboxsep
 \colorbox{shadecolor}{##1}\hskip-\fboxsep
     % There is no \\@totalrightmargin, so:
     \hskip-\linewidth \hskip-\@totalleftmargin \hskip\columnwidth}%
 \MakeFramed {\advance\hsize-\width
   \@totalleftmargin\z@ \linewidth\hsize
   \@setminipage}}%
 {\par\unskip\endMakeFramed%
 \at@end@of@kframe}
\makeatother

\definecolor{shadecolor}{rgb}{.97, .97, .97}
\definecolor{messagecolor}{rgb}{0, 0, 0}
\definecolor{warningcolor}{rgb}{1, 0, 1}
\definecolor{errorcolor}{rgb}{1, 0, 0}
\newenvironment{knitrout}{}{} % an empty environment to be redefined in TeX

\usepackage{alltt}

% Set lecture number for later use


% Part common to all the lectures
\subtitle{MATH 2740 -- Mathematics of Data Science -- Unit U09}
\author{\texorpdfstring{Julien Arino\newline\url{julien.arino@umanitoba.ca}}{Julien Arino}}
\institute{Department of Mathematics @ University of Manitoba}

% Title of the lecture
\title{Absorbing Markov chains, spectral centrality and PageRank}



\usetheme{default}
% Slide setup, colour independent

\usepackage{amsmath,amssymb,amsthm}
\usepackage[utf8]{inputenc}
\usepackage{colortbl}
\usepackage{bm}
\usepackage{xcolor}
\usepackage{dsfont}
\usepackage{setspace}
% To use \ding{234} and the like
\usepackage{pifont}
% To cross reference between slide files
\usepackage{zref-xr,zref-user}
\zxrsetup{tozreflabel=true,toltxlabel=true}
% Use something like
% \zexternaldocument{fileI}
% in the tex files. And cite using \zref instead of \ref

% Cross-reference system - see CROSS-REFERENCE-SETUP.md for manual setup instructions
\usepackage{booktabs}
\usepackage{marvosym}
\usepackage{cancel}
%\usepackage{transparent}
% Make doi clickable in the bibliography?
\usepackage{doi}

\usepackage[T1]{fontenc}

\usepackage{longtable}

% For heavier titles
\usepackage{helvet} % Enables Helvetica font family


% Fields and the like
\def\IC{\mathbb{C}}
\def\IE{\mathbb{E}}
\def\IF{\mathbb{F}}
\def\II{\mathbb{I}}
\def\IJ{\mathbb{J}}
\def\IK{\mathbb{K}}
\def\IM{\mathbb{M}}
\def\IN{\mathbb{N}}
\def\IP{\mathbb{P}}
\def\IR{\mathbb{R}}
\newcommand{\IRplus}{\mathbb{R}_+}
\def\IZ{\mathbb{Z}}
\def\11{\mathds{1}}


% Bold lowercase
\def\ba{\bm{a}}
\def\bb{\bm{b}}
\def\bc{\bm{c}}
\def\bd{\bm{d}}
\def\be{\bm{e}}
\def\bf{\bm{f}}
\def\bg{\bm{g}}
\def\bh{\bm{h}}
\def\bi{\bm{i}}
\def\bj{\bm{j}}
\def\bk{\bm{k}}
\def\bl{\bm{l}}
% Note: \bm is the command from the bm package, so \bm cannot be defined as \bm{m}
\def\bn{\bm{n}}
\def\bo{\bm{o}}
\def\bp{\bm{p}}
\def\bq{\bm{q}}
\def\br{\bm{r}}
\def\bs{\bm{s}}
\def\bt{\bm{t}}
\def\bu{\bm{u}}
\def\bv{\bm{v}}
\def\bw{\bm{w}}
\def\bx{\bm{x}}
\def\by{\bm{y}}
\def\bz{\bm{z}}
\newcommand{\vect}[1]{\bm{#1}}

% Bold capitals
\def\bA{\bm{A}}
\def\bB{\bm{B}}
\def\bC{\bm{C}}
\def\bD{\bm{D}}
\def\bE{\bm{E}}
\def\bF{\bm{F}}
\def\bG{\bm{G}}
\def\bH{\bm{H}}
\def\bI{\bm{I}}
\def\bJ{\bm{J}}
\def\bK{\bm{K}}
\def\bL{\bm{L}}
\def\bM{\bm{M}}
\def\bN{\bm{N}}
\def\bO{\bm{O}}
\def\bP{\bm{P}}
\def\bQ{\bm{Q}}
\def\bR{\bm{R}}
\def\bS{\bm{S}}
\def\bT{\bm{T}}
\def\bU{\bm{U}}
\def\bV{\bm{V}}
\def\bW{\bm{W}}
\def\bX{\bm{X}}
\def\bY{\bm{Y}}
\def\bZ{\bm{Z}}

% Bold numbers
\def\b0{\bm{0}}
\def\bzero{\bm{0}}
\def\bone{\bm{1}}

% Bold greek
\bmdefine{\bmu}{\bm{\mu}}
\def\bphi{\bm{\phi}}
\def\bvarphi{\bm{\varphi}}
\def\bPi{\bm{\Pi}}
\def\bGamma{\bm{\Gamma}}

% Bold red sentence
\def\boldred#1{{\color{red}\textbf{#1}}}
\def\defword#1{{\color{orange}\textbf{#1}}}

% Caligraphic letters
\def\A{\mathcal{A}}
\def\B{\mathcal{B}}
\def\C{\mathcal{C}}
\def\D{\mathcal{D}}
\def\E{\mathcal{E}}
\def\F{\mathcal{F}}
\def\G{\mathcal{G}}
\def\H{\mathcal{H}}
\def\I{\mathcal{I}}
\def\L{\mathcal{L}}
\def\M{\mathcal{M}}
\def\N{\mathcal{N}}
\def\P{\mathcal{P}}
\def\R{\mathcal{R}}
\def\S{\mathcal{S}}
\def\T{\mathcal{T}}
\def\U{\mathcal{U}}
\def\V{\mathcal{V}}

% Adding space for prime (') where needed
\def\pprime{\,'}
% Adding space for star (\star) where needed
\def\pstar{{\,\star}}

% tt font for code
\def\code#1{{\tt #1}}

% i.e., e.g.
\def\eg{\emph{e.g.}}
\def\ie{\emph{i.e.}}


% Operators and special symbols
\def\nbOne{{\mathchoice {\rm 1\mskip-4mu l} {\rm 1\mskip-4mu l}
{\rm 1\mskip-4.5mu l} {\rm 1\mskip-5mu l}}}
\def\cov{\ensuremath{\mathsf{cov}}}
\def\Var{\ensuremath{\mathsf{Var}\ }}
\def\Im{\textrm{Im}\;}
\def\Re{\textrm{Re}\;}
\def\det{\ensuremath{\mathsf{det}}}
\def\diag{\ensuremath{\mathsf{diag}}}
\def\nullspace{\ensuremath{\mathsf{null}}}
\def\nullity{\ensuremath{\mathsf{nullity}}}
\def\rank{\ensuremath{\mathsf{rank}}}
\def\range{\ensuremath{\mathsf{range}}}
\def\sgn{\ensuremath{\mathsf{sgn}}}
\def\Span{\ensuremath{\mathsf{span}}}
\def\tr{\ensuremath{\mathsf{tr}}}
\def\imply{$\Rightarrow$}
\def\restrictTo#1#2{\left.#1\right|_{#2}}
\newcommand{\parallelsum}{\mathbin{\!/\mkern-5mu/\!}}
\def\dsum{\mathop{\displaystyle \sum }}%
\def\dind#1#2{_{\substack{#1\\ #2}}}

\newcommand{\Qmatrix}[1]{%
  \begin{pmatrix}#1\end{pmatrix}%
}

\DeclareMathOperator{\GL}{GL}
\DeclareMathOperator{\Rel}{Re}
\def\Nt#1{\left|\!\left|\!\left|#1\right|\!\right|\!\right|}
\newcommand{\tripbar}{|\! |\! |}



% The beamer bullet (in base colour)
\def\bbullet{\leavevmode\usebeamertemplate{itemize item}\ }

% Colorblind-friendly colors (Okabe-Ito palette)
\definecolor{importantcolour}{RGB}{213,94,0}  % Vermilion (colorblind-safe highlight/alert)
\definecolor{othercolour}{RGB}{0,114,178}     % Deep/Sky Blue (colorblind-safe standard)
\colorlet{regularcolour}{othercolour}
\colorlet{standardcolour}{othercolour}

% Theorems and the like
\newtheorem{importantdefinition}[theorem]{Definition}
\newtheorem{importanttheorem}[theorem]{Theorem}
\newtheorem{importantproperty}[theorem]{Property}
\newtheorem{importantproposition}[theorem]{Proposition}
\newtheorem{importantcorollary}[theorem]{Corollary}
\newtheorem{importantlemma}[theorem]{Lemma}
\newtheorem{importantremark}[theorem]{Remark}
\newtheorem{proposition}[theorem]{Proposition}
\newtheorem{property}[theorem]{Property}
%\newtheorem{lemma}[theorem]{Lemma}
%\newtheorem{corollary}[theorem]{Corollary}
\newtheorem{remark}[theorem]{Remark}
\setbeamertemplate{theorems}[numbered]
%\setbeamertemplate{theorems}[ams style]

%
%\usecolortheme{orchid}
%\usecolortheme{orchid}

\def\red{\color[rgb]{1,0,0}}
\def\blue{\color[rgb]{0,0,1}}
\def\green{\color[rgb]{0,1,0}}
\def\important{\color{importantcolour}}
\def\boldimportant#1{{\color{importantcolour}\textbf{#1}}}

% Fix skipping lines after items in the bibliography
\setbeamertemplate{bibliography entry title}{}
\setbeamertemplate{bibliography entry location}{}
\setbeamertemplate{bibliography entry note}{}

% Get rid of navigation stuff
\setbeamertemplate{navigation symbols}{}

% Set footline/header line
\setbeamertemplate{footline}
{%
\quad p. \insertpagenumber \quad--\quad \insertsection\vskip2pt
}
% \setbeamertemplate{headline}
% {%
% \quad\insertsection\hfill p. \insertpagenumber\quad\mbox{}\vskip2pt
% }


\makeatletter
\newlength\beamerleftmargin
\setlength\beamerleftmargin{\Gm@lmargin}
\makeatother

% Colours for special pages
\def\extraContent{yellow!20}


%%%%%%%%%%%%%%%%%
\usepackage{tikz}
\usetikzlibrary{shapes,arrows}
\usetikzlibrary{positioning}
\usetikzlibrary{shapes.symbols,shapes.callouts,patterns}
\usetikzlibrary{calc,fit}
\usetikzlibrary{backgrounds}
\usetikzlibrary{decorations.pathmorphing,fit,petri}
\usetikzlibrary{automata}
\usetikzlibrary{fadings}
\usetikzlibrary{patterns,hobby}
\usetikzlibrary{backgrounds,fit,petri}
\usetikzlibrary{tikzmark}

\usepackage{pgfplots}
\pgfplotsset{compat=1.6}
\pgfplotsset{ticks=none}

\usetikzlibrary{decorations.markings}
\usetikzlibrary{arrows.meta}
\tikzset{>=stealth}

% For tikz
\tikzstyle{cloud} = [draw, ellipse,fill=red!20, node distance=0.87cm,
minimum height=2em]
\tikzstyle{line} = [draw, -latex']


%%% For max frame images
\newenvironment{changemargin}[2]{%
\begin{list}{}{%
\setlength{\topsep}{0pt}%
\setlength{\leftmargin}{#1}%
\setlength{\rightmargin}{#2}%
\setlength{\listparindent}{\parindent}%
\setlength{\itemindent}{\parindent}%
\setlength{\parsep}{\parskip}%
}%
\item[]}{\end{list}}


% Make one image take up the entire slide content area in beamer,.:
% centered/centred full-screen image, with title:
% This uses the whole screen except for the 1cm border around it
% all. 128x96mm
\newcommand{\titledFrameImage}[2]{
\begin{frame}{#1}
%\begin{changemargin}{-1cm}{-1cm}
\begin{center}
\includegraphics[width=108mm,height=\textheight,keepaspectratio]{#2}
\end{center}
%\end{changemargin}
\end{frame}
}

% Make one image take up the entire slide content area in beamer.:
% centered/centred full-screen image, no title:
% This uses the whole screen except for the 1cm border around it
% all. 128x96mm
\newcommand{\plainFrameImage}[1]{
\begin{frame}[plain]
%\begin{changemargin}{-1cm}{-1cm}
\begin{center}
\includegraphics[width=108mm,height=76mm,keepaspectratio]{#1}
\end{center}
%\end{changemargin}
\end{frame}
}

% Make one image take up the entire slide area, including borders, in beamer.:
% centered/centred full-screen image, no title:
% This uses the entire whole screen
\newcommand{\maxFrameImage}[1]{
\begin{frame}[plain]
\begin{changemargin}{-1cm}{-1cm}
\begin{center}
\includegraphics[width=\paperwidth,height=\paperheight,keepaspectratio]
{#1}
\end{center}
\end{changemargin}
\end{frame}
}

% This uses the entire whole screen (to include in frame)
\newcommand{\maxFrameImageNoFrame}[1]{
\begin{changemargin}{-1cm}{-1cm}
\begin{center}
\includegraphics[width=\paperwidth,height=0.99\paperheight,keepaspectratio]
{#1}
\end{center}
\end{changemargin}
}

% Make one image take up the entire slide area, including borders, in beamer.:
% centered/centred full-screen image, no title:
% This uses the entire whole screen
\newcommand{\maxFrameImageColor}[2]{
\begin{frame}[plain]
\setbeamercolor{normal text}{bg=#2!20}
\begin{changemargin}{-1cm}{-1cm}
\begin{center}
\includegraphics[width=\paperwidth,height=\paperheight,keepaspectratio]
{#1}
\end{center}
\end{changemargin}
\end{frame}
}


\usepackage{tikz}
\usetikzlibrary{patterns,hobby,matrix}
\usepackage{pgfplots}
\pgfplotsset{compat=1.6}
\pgfplotsset{ticks=none}

\usetikzlibrary{backgrounds}
\usetikzlibrary{decorations.markings}
\usetikzlibrary{arrows.meta}
\tikzset{>=stealth}

\tikzset{
  clockwise arrows/.style={
    postaction={
      decorate,
      decoration={
        markings,
        mark=between positions 0.1 and 0.9 step 40pt with {\arrow{>}},
   }}}}


% New integrated section command: creates section and section slide
\newcommand{\Ssection}[2]{
\section{#1}
\begin{frame}[noframenumbering,plain]
  \begin{tikzpicture}[remember picture,overlay]
    \node[above right,inner sep=0pt,opacity=0.2] at (current page.south west)
    {
        \includegraphics[height=\paperheight,width=\paperwidth]{#2}
    };
  \end{tikzpicture}
  \setbeamercolor{section in toc}{fg=section_page_list_colour}
  \setbeamerfont{section in toc}{size=\Large,series=\bfseries}
  \setbeamertemplate{section in toc shaded}[default][60]
  \tableofcontents[
    currentsection,
    sectionstyle=show/shaded,
    subsectionstyle=show/hide/hide,
    subsubsectionstyle=hide/hide/hide]
\end{frame}
\addtocounter{page}{-1}
}

% New integrated section command with subsections: creates section and section slide showing subsections
\newcommand{\SsectionWithSubs}[2]{
\section{#1}
\begin{frame}[noframenumbering,plain]
  \begin{tikzpicture}[remember picture,overlay]
    \node[above right,inner sep=0pt,opacity=0.2] at (current page.south west)
    {
        \includegraphics[height=\paperheight,width=\paperwidth]{#2}
    };
  \end{tikzpicture}
  \setbeamercolor{section in toc}{fg=section_page_list_colour}
  \setbeamerfont{section in toc}{size=\Large,series=\bfseries}
  \setbeamertemplate{section in toc shaded}[default][60]
  \tableofcontents[
    currentsection,
    sectionstyle=show/hide,
    subsectionstyle=show/show/hide,
    subsubsectionstyle=hide/hide/hide]
\end{frame}
\addtocounter{page}{-1}
}

% New integrated subsection command: creates subsection and subsection slide
\newcommand{\Ssubsection}[2]{
\subsection{#1}
\begin{frame}[noframenumbering,plain]
  \begin{tikzpicture}[remember picture,overlay]
    \node[above right,inner sep=0pt,opacity=0.2] at (current page.south west)
    {
        \includegraphics[height=\paperheight,width=\paperwidth]{#2}
    };
  \end{tikzpicture}
  \setbeamercolor{section in toc}{fg=subsection_page_list_colour}
  \setbeamerfont{section in toc}{size=\Large,series=\bfseries}
  \setbeamertemplate{section in toc shaded}[default][60]
  \setbeamerfont{subsection in toc}{series=\bfseries}
  \setbeamertemplate{subsection in toc shaded}[default][50]
  \tableofcontents[
    currentsection,
    sectionstyle=show/hide,
    subsectionstyle=show/shaded/hide,
    subsubsectionstyle=hide/hide/hide]
\end{frame}
\addtocounter{page}{-1}
}

% New integrated subsubsection command: creates subsubsection and subsubsection slide
\newcommand{\Ssubsubsection}[2]{
\subsubsection{#1}
\begin{frame}[noframenumbering,plain]
  \begin{tikzpicture}[remember picture,overlay]
    \node[above right,inner sep=0pt,opacity=0.2] at (current page.south west)
    {
        \includegraphics[height=\paperheight,width=\paperwidth]{#2}
    };
  \end{tikzpicture}
  \setbeamercolor{section in toc}{fg=subsub_header_section}
  \setbeamerfont{section in toc}{size=\Large,series=\bfseries}
  \setbeamertemplate{section in toc shaded}[default][60]
  \setbeamerfont{subsection in toc}{series=\bfseries}
  \setbeamertemplate{subsection in toc shaded}[default][50]
  \setbeamertemplate{subsubsection in toc shaded}[default][50]
  \tableofcontents[
    currentsection,
    sectionstyle=show/hide,
    subsectionstyle=show/hide/hide,
    subsubsectionstyle=show/shaded/hide]
\end{frame}
\addtocounter{page}{-1}
}

% Legacy commands (kept for backward compatibility)
% Beginning of a section
\newcommand{\newSectionSlide}[1]{
\begin{frame}[noframenumbering,plain]
  \begin{tikzpicture}[remember picture,overlay]
    \node[above right,inner sep=0pt,opacity=0.2] at (current page.south west)
    {
        \includegraphics[height=\paperheight,width=\paperwidth]{#1}
    };
  \end{tikzpicture}
  \setbeamercolor{section in toc}{fg=section_page_list_colour}
  \setbeamerfont{section in toc}{size=\Large,series=\bfseries}
  \setbeamertemplate{section in toc shaded}[default][60]
  \tableofcontents[
    currentsection,
    sectionstyle=show/shaded,
    subsectionstyle=show/hide/hide,
    subsubsectionstyle=hide/hide/hide]
\end{frame}
\addtocounter{page}{-1}
}

% Beginning of a section in which we also show subsections
\newcommand{\newSectionWithSubsSlide}[1]{
	\begin{frame}[noframenumbering,plain]
		\begin{tikzpicture}[remember picture,overlay]
			\node[above right,inner sep=0pt,opacity=0.2] at (current page.south west)
			{
				\includegraphics[height=\paperheight,width=\paperwidth]{#1}
			};
		\end{tikzpicture}
		\setbeamercolor{section in toc}{fg=section_page_list_colour}
		\setbeamerfont{section in toc}{size=\Large,series=\bfseries}
		\setbeamertemplate{section in toc shaded}[default][60]
		\tableofcontents[
		currentsection,
		sectionstyle=show/hide,
		subsectionstyle=show/show/hide,
		subsubsectionstyle=hide/hide/hide]
	\end{frame}
	\addtocounter{page}{-1}
}

% Beginning of a subsection
\newcommand{\newSubSectionSlide}[1]{
\begin{frame}[noframenumbering,plain]
  \begin{tikzpicture}[remember picture,overlay]
    \node[above right,inner sep=0pt,opacity=0.2] at (current page.south west)
    {
        \includegraphics[height=\paperheight,width=\paperwidth]{#1}
    };
  \end{tikzpicture}
  \setbeamercolor{section in toc}{fg=subsection_page_list_colour}
  \setbeamerfont{section in toc}{size=\Large,series=\bfseries}
  \setbeamertemplate{section in toc shaded}[default][60]
  \setbeamerfont{subsection in toc}{series=\bfseries}
  \setbeamertemplate{subsection in toc shaded}[default][50]
  \tableofcontents[
    currentsection,
    sectionstyle=show/hide,
    subsectionstyle=show/shaded/hide,
    subsubsectionstyle=hide/hide/hide]
\end{frame}
\addtocounter{page}{-1}
}

% Beginning of a subsubsection
\newcommand{\newSubSubSectionSlide}[1]{
\begin{frame}[noframenumbering,plain]
  \begin{tikzpicture}[remember picture,overlay]
    \node[above right,inner sep=0pt,opacity=0.2] at (current page.south west)
    {
        \includegraphics[height=\paperheight,width=\paperwidth]{#1}
    };
  \end{tikzpicture}
  \setbeamercolor{section in toc}{fg=subsub_header_section}
  \setbeamerfont{section in toc}{size=\Large,series=\bfseries}
  \setbeamertemplate{section in toc shaded}[default][60]
  \setbeamerfont{subsection in toc}{series=\bfseries}
  \setbeamertemplate{subsection in toc shaded}[default][50]
  \setbeamertemplate{subsubsection in toc shaded}[default][50]
  \tableofcontents[
    currentsection,
    sectionstyle=show/hide,
    subsectionstyle=show/hide/hide,
    subsubsectionstyle=show/shaded/hide]
\end{frame}
\addtocounter{page}{-1}
}


   %%%%%%%%%%%
% To have links to parts in the outline
\makeatletter
\AtBeginPart{%
  \addtocontents{toc}{\protect\beamer@partintoc{\the\c@part}{\beamer@partnameshort}{\the\c@page}}%
}
%% number, shortname, page.
\providecommand\beamer@partintoc[3]{%
  \ifnum\c@tocdepth=-1\relax
    % requesting onlyparts.
    \makebox[6em]{Part #1:} \textcolor{green!30!blue}{\hyperlink{#2}{#2}}
    \par
  \fi
}
\define@key{beamertoc}{onlyparts}[]{%
  \c@tocdepth=-1\relax
}
\makeatother%

\newcommand{\nameofthepart}{}
\newcommand{\nupart}[1]%
    {   \part{#1}%
        \renewcommand{\nameofthepart}{#1}%
        {
          \setbeamercolor{background canvas}{bg=orange!50}
          \begin{frame}{#1}%\partpage 
          \hypertarget{\nameofthepart}{}\tableofcontents%
          \end{frame}
        }
    }

% This command creates a title page using TikZ only
\newcommand{\tikztitlepage}[1]{%
\begin{frame}[plain,noframenumbering]
  \begin{tikzpicture}[remember picture,overlay]
    % Background image
    \node[above right,inner sep=0pt,opacity=0.1] 
      at (current page.south west) 
      {\includegraphics[width=\paperwidth,height=\paperheight]{#1}};

    % University logo
    \node[anchor=north east, inner sep=5pt, opacity=0.9] 
      at (current page.north east)
      {\includegraphics[width=0.2\textwidth]{FIGS-slides-admin/UM-logo-horizontal-CMYK.png}};
    
    % Title
    \node[anchor=center, align=center, 
          font=\fontsize{13}{15}\bfseries\color{UMbrown}, 
          text width=0.9\textwidth] 
          at ([yshift=2cm]current page.center)
          {\inserttitle};

      % Authors
      \node[anchor=center, align=center,
        font=\fontsize{10}{12}\bfseries\color{UMbrown},
        text width=0.7\textwidth]
        at ([yshift=0.8cm]current page.center)
        {\insertauthor};

      % Affiliation
      \node[anchor=north, align=center,
        font=\fontsize{9}{11}\color{UMbrown},
        text width=0.7\textwidth]
        at ([yshift=-0.2cm]current page.center)
        {\insertaffiliation};      
    % Date
    \node[anchor=north, align=center, 
          font=\fontsize{12}{16}\bfseries\color{UMbrown},
          text width=0.7\textwidth] 
          at ([yshift=0.2cm]current page.center)
          {\insertdate};

    % Land acknowledgement
    \node[anchor=south, align=justify, 
          font=\footnotesize, text=black, 
          text width=1.1\textwidth] 
          at ([yshift=0.5cm]current page.south)
          {The University of Manitoba campuses are located on original lands of Anishinaabeg, Ininew, Anisininew, Dakota and Dene peoples, and on the National Homeland of the Red River Métis.\\
          We respect the Treaties that were made on these territories, we acknowledge the harms and mistakes of the past, and we dedicate ourselves to move forward in partnership with Indigenous communities in a spirit of Reconciliation and collaboration.};
  \end{tikzpicture}
  \addtocounter{page}{-1}
\end{frame}
}
% The title page with figure
% \newcommand{\titlepagewithfigure}[1]{%
%   \begin{frame}[noframenumbering,plain]
%     \begin{tikzpicture}[remember picture,overlay]
%       \node[above right,inner sep=0pt,opacity=0.1] at (current page.south west)
%       {
%           \includegraphics[height=\paperheight,width=\paperwidth]{#1}
%       };
%       \node[anchor=north east,
%       inner sep=5pt,
%       opacity=0.9] at (current page.north east)
%       {
%           \includegraphics[width=0.2\textwidth]{FIGS-slides-admin/UM-logo-horizontal-CMYK.png}
%       };
%       \node[anchor=south, 
%       align=justify, 
%       text=black, 
%       text width=1.1\textwidth,
%       font=\footnotesize]  (land_acknowledgement)
%       at (current page.south) 
%       {The University of Manitoba campuses are located on original lands of Anishinaabeg, Ininew, Anisininew, Dakota and Dene peoples, and on the National Homeland of the Red River Métis.
%       We respect the Treaties that were made on these territories, we acknowledge the harms and mistakes of the past, and we dedicate ourselves to move forward in partnership with Indigenous communities in a spirit of Reconciliation and collaboration.};  
%       % \node[align=center, anchor=south,
%       % above=0.5cm of land_acknowledgement,
%       % text=black,
%       % font=\bfseries] {\insertdate};
%   \end{tikzpicture}
%   \setbeamercolor{title}{fg=title_page_title_colour}
%   \setbeamerfont{title}{size=\Large,series=\bfseries}
%   \setbeamercolor{author}{fg=title_page_author_colour}
%   \setbeamerfont{author}{size=\large,series=\bfseries}
%   \setbeamercolor{institute}{fg=title_page_institute_colour}
%   \setbeamerfont{institute}{size=\large,series=\bfseries}
%   \setbeamercolor{date}{fg=title_page_date_colour}
%   \setbeamerfont{date}{series=\bfseries}
% 	\titlepage
% \end{frame}
% \addtocounter{page}{-1}
% }

\newcommand{\titlepagewithfigure}[1]{%
  \begin{frame}[noframenumbering,plain]
    \begin{tikzpicture}[remember picture,overlay]
      \node[above right,inner sep=0pt,opacity=0.1] at (current page.south west)
      {
          \includegraphics[height=\paperheight,width=\paperwidth]{#1}
      };
      \node[anchor=north east,
      inner sep=5pt,
      opacity=0.9] at (current page.north east)
      {
          \includegraphics[width=0.2\textwidth]{FIGS-slides-admin/UM-logo-horizontal-CMYK.png}
      };
      \node[anchor=south, 
      align=justify, 
      text=black, 
      text width=1.1\textwidth,
      font=\footnotesize]  (land_acknowledgement)
      at (current page.south) 
      {The University of Manitoba campuses are located on original lands of Anishinaabeg, Ininew, Anisininew, Dakota and Dene peoples, and on the National Homeland of the Red River Métis.
      We respect the Treaties that were made on these territories, we acknowledge the harms and mistakes of the past, and we dedicate ourselves to move forward in partnership with Indigenous communities in a spirit of Reconciliation and collaboration.};  
      % \node[align=center, anchor=south,
      % above=0.5cm of land_acknowledgement,
      % text=black,
      % font=\bfseries] {\insertdate};
  \end{tikzpicture}
  \setbeamercolor{title}{fg=title_page_title_colour}
  \setbeamerfont{title}{size=\Large,series=\bfseries,family=\usefont{T1}{phv}{b}{n}}
  \setbeamercolor{author}{fg=title_page_author_colour}
  \setbeamerfont{author}{size=\large,series=\bfseries,family=\usefont{T1}{phv}{b}{n}}
  \setbeamercolor{institute}{fg=title_page_institute_colour}
  \setbeamerfont{institute}{size=\large,series=\bfseries,family=\usefont{T1}{phv}{b}{n}}
  \setbeamercolor{date}{fg=title_page_date_colour}
  \setbeamerfont{date}{series=\bfseries,family=\usefont{T1}{phv}{b}{n}}
	\titlepage
\end{frame}
\addtocounter{page}{-1}
}
% The outline page, with figure
% \newcommand{\outlinepage}[1]{%
% \begin{frame}[noframenumbering,plain]
%   \begin{tikzpicture}[remember picture,overlay]
%     \node[above right,inner sep=0pt,opacity=0.2] at (current page.south west)
%     {
%         \includegraphics[height=\paperheight,width=\paperwidth]{#1}
%     };
%   \end{tikzpicture}
%   \setbeamercolor{section in toc}{fg=outline_page_list_colour}
%   \setbeamerfont{section in toc}{size=\Large,series=\bfseries,family=\sffamily}
%   \frametitle{\textcolor{outline_page_title_colour}{\LARGE\bfseries Outline}}
%   \tableofcontents[hideallsubsections]
% \end{frame}
% \addtocounter{page}{-1}
% }
% The outline page, with figure
\newcommand{\outlinepage}[1]{%
\begin{frame}[noframenumbering,plain]
  \begin{tikzpicture}[remember picture,overlay]
    \node[above right,inner sep=0pt,opacity=0.2] at (current page.south west)
    {
        \includegraphics[height=\paperheight,width=\paperwidth]{#1}
    };
  \end{tikzpicture}
  \setbeamercolor{section in toc}{fg=outline_page_list_colour}
  % Use Helvetica Bold only for the outline slide TOC
  \setbeamerfont{section in toc}{size=\Large,family=\usefont{T1}{phv}{b}{n}}
  % Use Helvetica Bold for the outline title
  \frametitle{\textcolor{outline_page_title_colour}{\usefont{T1}{phv}{b}{n}\LARGE Outline}}
  \tableofcontents[hideallsubsections]
\end{frame}
\addtocounter{page}{-1}
}


%\let\oldsection\section
%\renewcommand{\section}[2]{\oldsection[#1]\newSectionSlide[#2]}


%%%%%%%%%%%%%%%%%%%%%
% CUSTOM SLIDE BACKGROUNDS
%%%%%%%%%%%%%%%%%%%%%
% Define custom background templates for different colors
\defbeamertemplate*{background canvas}{importantcolour}{%
  \color{importantcolour!15}\vrule width\paperwidth height\paperheight%
}
\defbeamertemplate*{background canvas}{blue}{%
  \color{blue!15}\vrule width\paperwidth height\paperheight%
}
\defbeamertemplate*{background canvas}{green}{%
  \color{green!15}\vrule width\paperwidth height\paperheight%
}
\defbeamertemplate*{background canvas}{red}{%
  \color{importantcolour!15}\vrule width\paperwidth height\paperheight%
}
\defbeamertemplate*{background canvas}{yellow}{%
  \color{yellow!20}\vrule width\paperwidth height\paperheight%
}
\defbeamertemplate*{background canvas}{purple}{%
  \color{purple!15}\vrule width\paperwidth height\paperheight%
}
\defbeamertemplate*{background canvas}{orange}{%
  \color{orange!20}\vrule width\paperwidth height\paperheight%
}

% Define keys for the different background options
\makeatletter
\define@key{beamerframe}{importantcolour}[true]{\setbeamertemplate{background canvas}[importantcolour]}
\define@key{beamerframe}{blue}[true]{\setbeamertemplate{background canvas}[blue]}
\define@key{beamerframe}{green}[true]{\setbeamertemplate{background canvas}[green]}
\define@key{beamerframe}{red}[true]{\setbeamertemplate{background canvas}[importantcolour]}
\define@key{beamerframe}{yellow}[true]{\setbeamertemplate{background canvas}[yellow]}
\define@key{beamerframe}{purple}[true]{\setbeamertemplate{background canvas}[purple]}
\define@key{beamerframe}{orange}[true]{\setbeamertemplate{background canvas}[orange]}
\makeatother

% Reset to normal background for all frames by default
\BeforeBeginEnvironment{frame}{\setbeamertemplate{background canvas}[default]}

\newcommand{\punnett}[2]{
    \begin{center}
    \renewcommand{\arraystretch}{1.5} % Add space to rows
    \begin{tabular}{|c c | c|c|}
        \multicolumn{2}{c}{} & \multicolumn{2}{c}{\textbf{Father}} \\
        \multicolumn{2}{c}{} & #1 \\ \hline
        #2
    \end{tabular}
    \end{center}
}

% Colour definitions for Punnett squares
\usepackage{multirow}
\usepackage{colortbl}
\colorlet{punnettorange}{orange!30}
\colorlet{punnettblack}{black!50}
\colorlet{punnetttortie}{orange!60!black}

\usetikzlibrary{positioning, decorations.pathreplacing, arrows.meta}

% --- TIKZ STYLES FOR CAT DIAGRAMS ---
\tikzset{
    % This is the main placeholder style for the cat images
    catnode/.style={
        draw, 
        rectangle, 
        rounded corners, 
        minimum height=1.5cm, 
        minimum width=2.5cm, 
        align=center, 
        font=\small\bfseries
    },
    % Allele styles (the X and Y at the top/side)
    allele/.style={font=\Large\bfseries},
    female/.style={allele, color=purple!80!black},
    male/.style={allele, color=blue!80!black},
    % --- Placeholder styles for different cats ---
    % YOU CAN REPLACE THE CONTENTS OF THESE NODES WITH YOUR IMAGES
    orange_cat/.style={
        catnode, 
        fill=orange!30, 
        text=black
    },
    black_cat/.style={
        catnode, 
        fill=black!70, 
        text=white
    },
    tortie_cat/.style={
        catnode, 
        fill=orange!50!black, % A mix for tortoiseshell
        text=white
    }
}

\usetikzlibrary{
    arrows.meta, % For nicer arrow heads (e.g., -Latex)
    positioning, % For relative node placement (e.g., above=of)
    automata     % For state diagrams, loops
}
% --- END TIKZ STYLES ---

% Common date for all slides
\date{Fall 2026}
% Counter for random transition slides generated by R
\newcounter{sssectionbg}
\setcounter{sssectionbg}{1}
\newcommand{\SSsection}[1]{
  \Ssection{#1}{FIGS/random_bg_\thesssectionbg.pdf}
  \stepcounter{sssectionbg}
}

% Random transition slide generated by R for the title page
\newcommand{\SStitlepage}{
  \titlepagewithfigure{FIGS/random_title_bg.pdf}
}

% Random transition slide generated by R for the outline page
\newcommand{\SSoutlinepage}{
  \outlinepage{FIGS/random_outline_bg.pdf}
}
\newcommand{\SSsubsection}[1]{
  \Ssubsection{#1}{FIGS/random_bg_\thesssectionbg.pdf}
  \stepcounter{sssectionbg}
}

\newcommand{\SSsubsubsection}[1]{
  \Ssubsubsection{#1}{FIGS/random_bg_\thesssectionbg.pdf}
  \stepcounter{sssectionbg}
}



\usecolortheme{orchid}
%% Listings
\usepackage{listings}
\definecolor{mygreen}{rgb}{0,0.6,0}
\definecolor{mygray}{rgb}{0.5,0.5,0.5}
\definecolor{mymauve}{rgb}{0.58,0,0.82}
\definecolor{mygold}{rgb}{1,0.843,0}
\definecolor{myblue}{rgb}{0.537,0.812,0.941}

\definecolor{mygold2}{RGB}{120,105,22}
\definecolor{mygrey2}{RGB}{50,50,50}

\definecolor{lgreen}{rgb}{0.6,0.9,.6}
\definecolor{lred}{rgb}{1,0.5,.5}

\lstloadlanguages{R}
\lstset{ %
  language=R,
  backgroundcolor=\color{black!05},   % choose the background color
  basicstyle=\footnotesize\ttfamily,        % size of fonts used for the code
  breaklines=true,                 % automatic line breaking only at whitespace
  captionpos=b,                    % sets the caption-position to bottom
  commentstyle=\color{mygreen},    % comment style
  escapeinside={\%*}{*)},          % if you want to add LaTeX within your code
  keywordstyle=\color{red},       % keyword style
  stringstyle=\color{mygold},     % string literal style
  keepspaces=true,
  columns=fullflexible,
  tabsize=4,
}
% Could also do (in lstset)
% basicstyle==\fontfamily{pcr}\footnotesize
\lstdefinelanguage{Renhanced}%
  {keywords={abbreviate,abline,abs,acos,acosh,action,add1,add,%
      aggregate,alias,Alias,alist,all,anova,any,aov,aperm,append,apply,%
      approx,approxfun,apropos,Arg,args,array,arrows,as,asin,asinh,%
      atan,atan2,atanh,attach,attr,attributes,autoload,autoloader,ave,%
      axis,backsolve,barplot,basename,besselI,besselJ,besselK,besselY,%
      beta,binomial,body,box,boxplot,break,browser,bug,builtins,bxp,by,%
      c,C,call,Call,case,cat,category,cbind,ceiling,character,char,%
      charmatch,check,chol,chol2inv,choose,chull,class,close,cm,codes,%
      coef,coefficients,co,col,colnames,colors,colours,commandArgs,%
      comment,complete,complex,conflicts,Conj,contents,contour,%
      contrasts,contr,control,helmert,contrib,convolve,cooks,coords,%
      distance,coplot,cor,cos,cosh,count,fields,cov,covratio,wt,CRAN,%
      create,crossprod,cummax,cummin,cumprod,cumsum,curve,cut,cycle,D,%
      data,dataentry,date,dbeta,dbinom,dcauchy,dchisq,de,debug,%
      debugger,Defunct,default,delay,delete,deltat,demo,de,density,%
      deparse,dependencies,Deprecated,deriv,description,detach,%
      dev2bitmap,dev,cur,deviance,off,prev,,dexp,df,dfbetas,dffits,%
      dgamma,dgeom,dget,dhyper,diag,diff,digamma,dim,dimnames,dir,%
      dirname,dlnorm,dlogis,dnbinom,dnchisq,dnorm,do,dotplot,double,%
      download,dpois,dput,drop,drop1,dsignrank,dt,dummy,dump,dunif,%
      duplicated,dweibull,dwilcox,dyn,edit,eff,effects,eigen,else,%
      emacs,end,environment,env,erase,eval,equal,evalq,example,exists,%
      exit,exp,expand,expression,External,extract,extractAIC,factor,%
      fail,family,fft,file,filled,find,fitted,fivenum,fix,floor,for,%
      For,formals,format,formatC,formula,Fortran,forwardsolve,frame,%
      frequency,ftable,ftable2table,function,gamma,Gamma,gammaCody,%
      gaussian,gc,gcinfo,gctorture,get,getenv,geterrmessage,getOption,%
      getwd,gl,glm,globalenv,gnome,GNOME,graphics,gray,grep,grey,grid,%
      gsub,hasTsp,hat,heat,help,hist,home,hsv,httpclient,I,identify,if,%
      ifelse,Im,image,\%in\%,index,influence,measures,inherits,install,%
      installed,integer,interaction,interactive,Internal,intersect,%
      inverse,invisible,IQR,is,jitter,kappa,kronecker,labels,lapply,%
      layout,lbeta,lchoose,lcm,legend,length,levels,lgamma,library,%
      licence,license,lines,list,lm,load,local,locator,log,log10,log1p,%
      log2,logical,loglin,lower,lowess,ls,lsfit,lsf,ls,machine,Machine,%
      mad,mahalanobis,make,link,margin,match,Math,matlines,mat,matplot,%
      matpoints,matrix,max,mean,median,memory,menu,merge,methods,min,%
      missing,Mod,mode,model,response,mosaicplot,mtext,mvfft,na,nan,%
      names,omit,nargs,nchar,ncol,NCOL,new,next,NextMethod,nextn,%
      nlevels,nlm,noquote,NotYetImplemented,NotYetUsed,nrow,NROW,null,%
      numeric,\%o\%,objects,offset,old,on,Ops,optim,optimise,optimize,%
      options,or,order,ordered,outer,package,packages,page,pairlist,%
      pairs,palette,panel,par,parent,parse,paste,path,pbeta,pbinom,%
      pcauchy,pchisq,pentagamma,persp,pexp,pf,pgamma,pgeom,phyper,pico,%
      pictex,piechart,Platform,plnorm,plogis,plot,pmatch,pmax,pmin,%
      pnbinom,pnchisq,pnorm,points,poisson,poly,polygon,polyroot,pos,%
      postscript,power,ppoints,ppois,predict,preplot,pretty,Primitive,%
      print,prmatrix,proc,prod,profile,proj,prompt,prop,provide,%
      psignrank,ps,pt,ptukey,punif,pweibull,pwilcox,q,qbeta,qbinom,%
      qcauchy,qchisq,qexp,qf,qgamma,qgeom,qhyper,qlnorm,qlogis,qnbinom,%
      qnchisq,qnorm,qpois,qqline,qqnorm,qqplot,qr,Q,qty,qy,qsignrank,%
      qt,qtukey,quantile,quasi,quit,qunif,quote,qweibull,qwilcox,%
      rainbow,range,rank,rbeta,rbind,rbinom,rcauchy,rchisq,Re,read,csv,%
      csv2,fwf,readline,socket,real,Recall,rect,reformulate,regexpr,%
      relevel,remove,rep,repeat,replace,replications,report,require,%
      resid,residuals,restart,return,rev,rexp,rf,rgamma,rgb,rgeom,R,%
      rhyper,rle,rlnorm,rlogis,rm,rnbinom,RNGkind,rnorm,round,row,%
      rownames,rowsum,rpois,rsignrank,rstandard,rstudent,rt,rug,runif,%
      rweibull,rwilcox,sample,sapply,save,scale,scan,scan,screen,sd,se,%
      search,searchpaths,segments,seq,sequence,setdiff,setequal,set,%
      setwd,show,sign,signif,sin,single,sinh,sink,solve,sort,source,%
      spline,splinefun,split,sqrt,stars,start,stat,stem,step,stop,%
      storage,strstrheight,stripplot,strsplit,structure,strwidth,sub,%
      subset,substitute,substr,substring,sum,summary,sunflowerplot,svd,%
      sweep,switch,symbol,symbols,symnum,sys,status,system,t,table,%
      tabulate,tan,tanh,tapply,tempfile,terms,terrain,tetragamma,text,%
      time,title,topo,trace,traceback,transform,tri,trigamma,trunc,try,%
      ts,tsp,typeof,unclass,undebug,undoc,union,unique,uniroot,unix,%
      unlink,unlist,unname,untrace,update,upper,url,UseMethod,var,%
      variable,vector,Version,vi,warning,warnings,weighted,weights,%
      which,while,window,write,\%x\%,x11,X11,xedit,xemacs,xinch,xor,%
      xpdrows,xy,xyinch,yinch,zapsmall,zip},%
   otherkeywords={!,!=,~,$,*,\%,\&,\%/\%,\%*\%,\%\%,<-,<<-,_,/},%
   alsoother={._$},%
   sensitive,%
   morecomment=[l]\#,%
   morestring=[d]",%
   morestring=[d]'% 2001 Robert Denham
  }%

% Custom color for listings
\definecolor{codegreen}{rgb}{0,0.6,0}
\definecolor{codegray}{rgb}{0.5,0.5,0.5}
\definecolor{codepurple}{rgb}{0.58,0,0.82}
\definecolor{backcolour}{rgb}{0.95,0.95,0.92}

%------------------------------------------------------------
% Listings configuration for Julia
%------------------------------------------------------------
% Define the Julia language for the listings package
\lstdefinelanguage{Julia}
{
  morekeywords={abstract,baremodule,begin,break,catch,const,continue,do,else,elseif,end,export,false,finally,for,function,global,if,import,let,local,macro,module,quote,return,struct,true,try,using,while,in},
  sensitive=true,
  comment=[l]{\#},
  morecomment=[n]{\#=}{=\#}, % For block comments like #= ... =#
  string=[b]",
  morestring=[b]',
  alsoletter={!}, % Allow ! in identifiers for mutating functions
}

% Define the style for the code blocks
\lstdefinestyle{JuliaStyle}{
    backgroundcolor=\color{backcolour},   
    commentstyle=\color{codegreen},
    keywordstyle=\color{magenta},
    numberstyle=\tiny\color{codegray},
    stringstyle=\color{codepurple},
    basicstyle=\ttfamily\footnotesize,
    breakatwhitespace=false,         
    breaklines=true,                 
    captionpos=b,                    
    keepspaces=true,                 
    % numbers=left,                    
    % numbersep=5pt,                  
    showspaces=false,                
    showstringspaces=false,
    showtabs=false,                  
    tabsize=2,
    language=Julia % Use the newly defined Julia language
}
\lstset{style=JuliaStyle}

%%%%%%% 
%% Definitions in yellow boxes
\usepackage{etoolbox}
\setbeamercolor{block title}{use=structure,fg=structure.fg,bg=structure.fg!40!bg}
\setbeamercolor{block body}{parent=normal text,use=block title,bg=block title.bg!20!bg}

% Important definition
\BeforeBeginEnvironment{importantdefinition}{%
	\setbeamercolor{block title}{fg=black,bg=yellow!20!white}
	\setbeamercolor{block body}{fg=black, bg=yellow!05!white}
}
\AfterEndEnvironment{importantdefinition}{
	\setbeamercolor{block title}{use=structure,fg=structure.fg,bg=structure.fg!20!bg}
	\setbeamercolor{block body}{parent=normal text,use=block title,bg=block title.bg!50!bg, fg=black}
}
% Regular (less important) definition
\BeforeBeginEnvironment{definition}{%
	\setbeamercolor{block title}{fg=black,bg=black!15!white}
	\setbeamercolor{block body}{fg=black, bg=black!05!white}
}
\AfterEndEnvironment{definition}{
	\setbeamercolor{block title}{use=structure,fg=structure.fg,bg=structure.fg!20!bg}
	\setbeamercolor{block body}{parent=normal text,use=block title,bg=block title.bg!50!bg, fg=black}
}
% Important theorem (red)
\BeforeBeginEnvironment{importanttheorem}{%
	\setbeamercolor{block title}{fg=black,bg=red!20!white}
	\setbeamercolor{block body}{fg=black, bg=red!05!white}
}
\AfterEndEnvironment{importanttheorem}{
	\setbeamercolor{block title}{use=structure,fg=structure.fg,bg=structure.fg!20!bg}
	\setbeamercolor{block body}{parent=normal text,use=block title,bg=block title.bg!50!bg, fg=black}
}
\BeforeBeginEnvironment{importantproperty}{%
	\setbeamercolor{block title}{fg=black,bg=red!50!white}
	\setbeamercolor{block body}{fg=black, bg=red!30!white}
}
\AfterEndEnvironment{importantproperty}{
	\setbeamercolor{block title}{use=structure,fg=structure.fg,bg=structure.fg!20!bg}
	\setbeamercolor{block body}{parent=normal text,use=block title,bg=block title.bg!50!bg, fg=black}
}

%% Colours for title page, outline and the like
\definecolor{UMbrown}{RGB}{79,44,29}
\definecolor{HighlightGold}{RGB}{230,165,83}
\definecolor{DeepGreen}{RGB}{27,46,28}
\definecolor{MediumGreen}{RGB}{52,128,56}
\definecolor{LightGreen}{RGB}{150,224,154}
\definecolor{AccentRed}{RGB}{196,44,27}
%% These are the titles used in practice
% Title page
\colorlet{title_page_title_colour}{UMbrown}
\colorlet{title_page_author_colour}{UMbrown}
\colorlet{title_page_date_colour}{UMbrown}
\colorlet{title_page_institute_colour}{UMbrown}
% Outline page
\colorlet{outline_page_title_colour}{DeepGreen}
\colorlet{outline_page_list_colour}{UMbrown}
% Section page
\colorlet{section_page_title_colour}{DeepGreen}
\colorlet{section_page_list_colour}{UMbrown}
% Subection page
\colorlet{subsection_page_title_colour}{DeepGreen}
\colorlet{subsection_page_list_colour}{UMbrown}



\providecommand{\bt}{\bm{t}}



\IfFileExists{upquote.sty}{\usepackage{upquote}}{}
\begin{document}

% Set up cross-references and counter persistence
\setcounter{theorem}{137}

%%%%%%%%%%%%%%%%%%%%%%%%%%%%%%%%%
%%%%%%%%%%%%%%%%%%%%%%%%%%%%%%%%%
%% TITLE AND OUTLINE
%%%%%%%%%%%%%%%%%%%%%%%%%%%%%%%%%
%%%%%%%%%%%%%%%%%%%%%%%%%%%%%%%%%
\SStitlepage
\SSoutlinepage

\begin{frame}{From Unit 08 to Unit 09}
In Unit 08, we studied irreducible and primitive nonnegative matrices and saw that a regular Markov chain converges to a unique positive stationary distribution $\bv$ satisfying $\bv^T P = \bv^T$ (or $P^T\bv = \bv$)
\vfill
Two natural questions arise when a directed network is \emph{not} regular:
\begin{enumerate}
\item What happens when a Markov chain has \defword{absorbing states} (states that can be entered but never left), making $P$ reducible?
\item How do we rank the importance of vertices in large directed networks (citations, social networks or the World Wide Web) using eigenvectors and linear systems?
\end{enumerate}
\vfill
Both questions are answered by the same two linear algebra tools: \textbf{Perron-Frobenius theory} and the \textbf{Neumann series} $(\II - M)^{-1} = \sum_{k=0}^\infty M^k$
\end{frame}


%%%%%%%%%%%%%%%%%%%%%%%%%%%%%%%%%
%%%%%%%%%%%%%%%%%%%%%%%%%%%%%%%%%
%% SECTION 1: ABSORBING MARKOV CHAINS
%%%%%%%%%%%%%%%%%%%%%%%%%%%%%%%%%
%%%%%%%%%%%%%%%%%%%%%%%%%%%%%%%%%
\SSsection{Absorbing Markov chains}

\SSsubsection{Absorbing states and motivating examples}

\begin{frame}{Motivating context: the gambler's ruin}
Consider the classical \defword{gambler's ruin} problem: a gambler starts with $\$k$ ($k \ge 1$) and repeatedly bets $\$1$ against an infinitely rich house
\begin{itemize}
\item On each bet, the gambler wins $\$1$ with probability $p \in (0, 1)$ and loses $\$1$ with probability $q = 1 - p$
\item If the gambler's fortune reaches $\$0$, play stops forever: state $0$ is the \textbf{sole absorbing state}
\end{itemize}
\vfill
\begin{center}
\begin{tikzpicture}[
    ->,
    >=Latex,
    node distance=1.8cm,
    every path/.style={thick},
    abs/.style={state, fill=red!20, minimum size=0.95cm},
    trans/.style={state, fill=black!10, minimum size=0.95cm}
]
    \node[abs] (S0) {$0$};
    \node[trans, right=of S0] (S1) {$1$};
    \node[trans, right=of S1] (S2) {$2$};
    \node[trans, right=of S2] (S3) {$3$};
    \node[right=1.2cm of S3] (dots) {$\cdots$};
    \path (S0) edge [loop left] node[left] {$1$} (S0);
    \path (S1) edge node[above] {$q$} (S0);
    \path (S1) edge [bend left=18] node[above] {$p$} (S2);
    \path (S2) edge [bend left=18] node[below] {$q$} (S1);
    \path (S2) edge [bend left=18] node[above] {$p$} (S3);
    \path (S3) edge [bend left=18] node[below] {$q$} (S2);
    \path (S3) edge [bend left=18] node[above] {$p$} (dots);
    \path (dots) edge [bend left=18] node[below] {$q$} (S3);
\end{tikzpicture}
\end{center}
\vfill
If the gambler also stops upon reaching a target jackpot, we obtain a finite chain with \textbf{two} absorbing boundaries
\end{frame}

\begin{frame}{Motivating example 1: 1D random walk with absorbing walls}
Consider a particle (or gambler's fortune) moving on the states $\{S_1, S_2, S_3, S_4, S_5\}$:
\begin{itemize}
\item From an interior state $S_i$ ($i \in \{2,3,4\}$), the particle moves left to $S_{i-1}$ with probability $q = 1-p$ and right to $S_{i+1}$ with probability $p \in (0,1)$
\item States $S_1$ (ruin) and $S_5$ (target jackpot) are \textbf{absorbing boundaries}: once reached, the process stays there forever
\end{itemize}
\vfill
\begin{center}
\begin{tikzpicture}[
    ->,
    >=Latex,
    node distance=1.6cm,
    every path/.style={thick},
    abs/.style={state, fill=red!20, minimum size=1cm},
    trans/.style={state, fill=black!10, minimum size=1cm}
]
    \node[abs] (S1) {$S_1$};
    \node[trans, right=of S1] (S2) {$S_2$};
    \node[trans, right=of S2] (S3) {$S_3$};
    \node[trans, right=of S3] (S4) {$S_4$};
    \node[abs, right=of S4] (S5) {$S_5$};
    \path (S1) edge [loop left] node[left] {$1$} (S1);
    \path (S5) edge [loop right] node[right] {$1$} (S5);
    \path (S2) edge node[above] {$q$} (S1);
    \path (S2) edge [bend left=18] node[above] {$p$} (S3);
    \path (S3) edge [bend left=18] node[below] {$q$} (S2);
    \path (S3) edge [bend left=18] node[above] {$p$} (S4);
    \path (S4) edge [bend left=18] node[below] {$q$} (S3);
    \path (S4) edge node[above] {$p$} (S5);
\end{tikzpicture}
\end{center}
\end{frame}

\begin{frame}{Motivating example 2: modifying the breeding experiment}
We modify the sex-linked orange cat breeding setup from Unit 08 by repeatedly mating a female offspring with an \textbf{orange male} ($X^O Y$) at every generation
\vfill
The state of the chain is the genotype of the female offspring:
\begin{itemize}
\item $S_1$: $X^O X^O$ (orange female)
\item $S_2$: $X^o X^o$ (black female)
\item $S_3$: $X^O X^o$ (tortoiseshell female)
\end{itemize}
\vfill
Let us compute the transition probabilities $p_{ij} = \IP(S_j(t+1) \mid S_i(t))$ when mating each female genotype with an $X^O Y$ male
\end{frame}

\begin{frame}{Transitions for the modified breeding experiment}
\begin{columns}[T]
\begin{column}{0.32\textwidth}
\centering
\textbf{Current state $S_1$ ($X^O X^O$)}
\vspace{0.2cm}

\renewcommand{\arraystretch}{1.4}
\begin{tabular}{c c | c | c |}
\multicolumn{2}{c}{} & \multicolumn{1}{c}{$X^O$} & \multicolumn{1}{c}{$Y$} \\ \cline{3-4}
\multirow{2}{*}{\rotatebox{90}{Mom}} & $X^O$ & \cellcolor{punnettorange}$S_1$ & \cellcolor{punnettorange}M \\ \cline{2-4}
& $X^O$ & \cellcolor{punnettorange}$S_1$ & \cellcolor{punnettorange}M \\ \cline{2-4}
\end{tabular}
\vspace{0.2cm}

100\% $X^O X^O$ ($S_1$)\\
$\implies p_{11} = 1$
\end{column}
\begin{column}{0.32\textwidth}
\centering
\textbf{Current state $S_2$ ($X^o X^o$)}
\vspace{0.2cm}

\renewcommand{\arraystretch}{1.4}
\begin{tabular}{c c | c | c |}
\multicolumn{2}{c}{} & \multicolumn{1}{c}{$X^O$} & \multicolumn{1}{c}{$Y$} \\ \cline{3-4}
\multirow{2}{*}{\rotatebox{90}{Mom}} & $X^o$ & \cellcolor{punnetttortie}$S_3$ & \cellcolor{punnettblack}M \\ \cline{2-4}
& $X^o$ & \cellcolor{punnetttortie}$S_3$ & \cellcolor{punnettblack}M \\ \cline{2-4}
\end{tabular}
\vspace{0.2cm}

100\% $X^O X^o$ ($S_3$)\\
$\implies p_{23} = 1$
\end{column}
\begin{column}{0.32\textwidth}
\centering
\textbf{Current state $S_3$ ($X^O X^o$)}
\vspace{0.2cm}

\renewcommand{\arraystretch}{1.4}
\begin{tabular}{c c | c | c |}
\multicolumn{2}{c}{} & \multicolumn{1}{c}{$X^O$} & \multicolumn{1}{c}{$Y$} \\ \cline{3-4}
\multirow{2}{*}{\rotatebox{90}{Mom}} & $X^O$ & \cellcolor{punnettorange}$S_1$ & \cellcolor{punnettorange}M \\ \cline{2-4}
& $X^o$ & \cellcolor{punnetttortie}$S_3$ & \cellcolor{punnettblack}M \\ \cline{2-4}
\end{tabular}
\vspace{0.2cm}

50\% $S_1$, 50\% $S_3$\\
$\implies p_{31} = 1/2,\; p_{33} = 1/2$
\end{column}
\end{columns}
\end{frame}

\begin{frame}{Summary of the modified breeding chain}
In our row-stochastic convention ($p_{ij} = \IP(S_j(t+1) \mid S_i(t))$), the transition matrix is
\[
P = \begin{pmatrix}
1 & 0 & 0 \\
0 & 0 & 1 \\
1/2 & 0 & 1/2
\end{pmatrix}
\]
\vfill
\begin{center}
\begin{tikzpicture}[
    ->,
    >=Latex,
    node distance=2.5cm,
    every path/.style={thick}
]
    \node[state,fill=punnettorange] (S1) at (0, 1.8) {$X^OX^O$};
    \node[state,fill=punnettblack,text=white] (S2) at (-2.5, 0) {$X^oX^o$};
    \node[state,fill=punnetttortie,text=white] (S3) at (2.5, 0) {$X^OX^o$};
    \path (S1) edge [loop left] node[left] {$1$} (S1);
    \path (S2) edge node[above] {$1$} (S3);
    \path (S3) edge [loop right] node[right] {$0.5$} (S3);
    \path (S3) edge node[sloped,above] {$0.5$} (S1);
\end{tikzpicture}
\end{center}
\end{frame}

\begin{frame}{Absorbing states}
\begin{importantdefinition}[Absorbing state]
A state $S_i$ in a Markov chain with row-stochastic transition matrix $P = [p_{ij}]$ is \defword{absorbing} if once the chain enters $S_i$, it never leaves:
\[
p_{ii} = 1 \quad \text{and} \quad p_{ij} = 0 \text{ for all } j \neq i
\]
Equivalently, row $i$ of $P$ is the standard basis row vector $\be_i^T$, and in $\G(P)$ vertex $S_i$ has a self-loop of weight $1$ and no outgoing arcs to any other vertex
\end{importantdefinition}
\end{frame}

\begin{frame}{Absorbing Markov chains}
\begin{importantdefinition}[Absorbing Markov chain]
A Markov chain is \defword{absorbing} if:
\begin{enumerate}
\item It has at least one absorbing state
\item From every non-absorbing state, it is possible to reach at least one absorbing state in a finite number of steps (there is a directed walk in $\G(P)$ to an absorbing state)
\end{enumerate}
In an absorbing Markov chain, a state that is not absorbing is called \defword{transient}
\end{importantdefinition}
\end{frame}

\begin{frame}{Key questions about absorbing Markov chains}
Suppose we have an absorbing chain such as the 5-state random walk with two absorbing boundaries $S_1$ and $S_5$ (and three transient states $S_2, S_3, S_4$):
\begin{center}
\begin{tikzpicture}[
    ->,
    >=Latex,
    node distance=1.8cm,
    every path/.style={thick},
    abs/.style={state, fill=red!20, minimum size=0.95cm},
    trans/.style={state, fill=black!10, minimum size=0.95cm}
]
    \node[abs] (1) {$S_1$};
    \node[trans, right=of 1] (2) {$S_2$};
    \node[trans, right=of 2] (3) {$S_3$};
    \node[trans, right=of 3] (4) {$S_4$};
    \node[abs, right=of 4] (5) {$S_5$};
    \draw[->] (1) to [loop left] (1);
    \draw[->] (5) to [loop right] (5);
    \draw[->] (2) to (1);
    \draw[->] (2) to [bend left=18] (3);
    \draw[->] (3) to [bend left=18] (2);
    \draw[->] (3) to [bend left=18] (4);
    \draw[->] (4) to [bend left=18] (3);
    \draw[->] (4) to (5);
\end{tikzpicture}
\end{center}
\vfill
\begin{enumerate}
\item Does the process eventually reach an absorbing state?
\item Starting in transient state $S_i$, what is the expected number of visits to transient state $S_j$?
\item Starting in transient state $S_i$, what is the expected number of steps before absorption?
\item When there are multiple absorbing states, what is the probability of being absorbed in absorbing state $S_j$ when starting in transient state $S_i$?
\end{enumerate}
\end{frame}


%%%%%%%%%%%%%%%%%%%%%%%%%%%%%%%%%
\SSsubsection{Canonical block form and reducibility}

\begin{frame}{Absorbing chains are reducible}
Suppose a Markov chain on $n$ states has $r \ge 1$ absorbing states and $m = n - r \ge 1$ transient states
\vfill
\begin{itemize}
\item Because an absorbing state $S_i$ has no outgoing arcs to any other vertex ($p_{ij} = 0$ for all $j \neq i$), there is no directed path from $S_i$ to any transient state
\item Therefore, the transition digraph $\G(P)$ is \textbf{not strongly connected}
\item By the characterisation from Unit 08, the transition matrix $P$ is \defword{reducible}
\end{itemize}
\vfill
Relabelling the states of the chain corresponds to applying a permutation matrix $\Pi \in \M_n$:
\[
\tilde{P} = \Pi^T P \Pi
\]
We always relabel the states so that the $r$ \textbf{absorbing states come first} ($S_1, \ldots, S_r$) and the $m = n - r$ \textbf{transient states come last} ($S_{r+1}, \ldots, S_n$)
\end{frame}

\begin{frame}{Example: 5-state random walk in original ordering}
In Motivating example 1 on $(S_1, S_2, S_3, S_4, S_5)$, states $S_1$ and $S_5$ are absorbing while $S_2, S_3, S_4$ are transient
\vfill
In the original state ordering, the transition matrix is
\[
P_{\text{orig}} = \begin{array}{c@{\hspace{4pt}}c}
& \begin{array}{ccccc} S_1 & S_2 & S_3 & S_4 & S_5 \end{array} \\
\begin{array}{c} S_1 \\ S_2 \\ S_3 \\ S_4 \\ S_5 \end{array} &
\left(\begin{array}{ccccc}
1 & 0 & 0 & 0 & 0 \\
q & 0 & p & 0 & 0 \\
0 & q & 0 & p & 0 \\
0 & 0 & q & 0 & p \\
0 & 0 & 0 & 0 & 1
\end{array}\right)
\end{array}
\]
\end{frame}

\begin{frame}{Example: reordering the 5-state random walk into canonical form}
Reordering the states as $(S_1, S_5, S_2, S_3, S_4)$ via the permutation matrix $\Pi = [\be_1\; \be_5\; \be_2\; \be_3\; \be_4]$ gives the canonical block form $\tilde{P} = \Pi^T P_{\text{orig}} \Pi$:
\vfill
\[
\tilde{P} = \begin{array}{c@{\hspace{4pt}}c}
& \begin{array}{cc|ccc} S_1 & S_5 & S_2 & S_3 & S_4 \end{array} \\
\begin{array}{c} S_1 \\ S_5 \\ \hline S_2 \\ S_3 \\ S_4 \end{array} &
\left(\begin{array}{cc|ccc}
1 & 0 & 0 & 0 & 0 \\
0 & 1 & 0 & 0 & 0 \\ \hline
q & 0 & 0 & p & 0 \\
0 & 0 & q & 0 & p \\
0 & p & 0 & q & 0
\end{array}\right)
\end{array} = \begin{pmatrix} \II_2 & \b0_{2,3} \\ R & Q \end{pmatrix}
\]
\end{frame}

\begin{frame}{Canonical block form of an absorbing chain}
\begin{importantdefinition}[Canonical block form]
For an absorbing Markov chain with $r$ absorbing states (listed first) and $m = n - r$ transient states (listed last), the row-stochastic transition matrix takes the \defword{canonical block form} (or \defword{standard form})
\[
P = \begin{pmatrix}
\II_r & \b0_{rm} \\
R & Q
\end{pmatrix}
\]
\end{importantdefinition}
\vfill
Here:
\begin{itemize}
\item $\II_r \in \M_r$ is the $r \times r$ identity matrix (absorbing states stay put)
\item $\b0_{rm} \in \M_{rm}$ is the $r \times m$ zero matrix (absorbing states never move to transient states)
\item $R \in \M_{mr}$ contains the one-step transition probabilities from transient to absorbing states
\item $Q \in \M_m$ contains the one-step transition probabilities among transient states
\end{itemize}
\end{frame}

\begin{frame}{Visualising the canonical blocks}
In our row-stochastic convention, rows correspond to the \textbf{from} state (at time $t$) and columns correspond to the \textbf{to} state (at time $t+1$):
\vfill
\begin{center}
\renewcommand{\arraystretch}{1.5}
\begin{tabular}{c|cc}
& \textbf{To absorbing} ($r$ cols) & \textbf{To transient} ($m$ cols) \\ \hline
\textbf{From absorbing} ($r$ rows) & $\II_r \in \M_r$ & $\b0_{rm} \in \M_{rm}$ \\
\textbf{From transient} ($m$ rows) & $R \in \M_{mr}$ & $Q \in \M_m$
\end{tabular}
\end{center}
\vfill
Notice that:
\begin{itemize}
\item Each row of $\begin{pmatrix} R & Q \end{pmatrix}$ sums to $1$: $R\11_r + Q\11_m = \11_m$
\item $Q$ is nonnegative and \textbf{row-substochastic}: $Q\11_m \le \11_m$, with strict inequality in every row $i$ that has a direct transition to an absorbing state ($\sum_{j=1}^r r_{ij} > 0$)
\end{itemize}
\end{frame}

\begin{frame}{Powers of $P$ in canonical form}
To understand the long-term behaviour $\bp(k)^T = \bp(0)^T P^k$, we compute the $k$th power of $P$ in block form
\vfill
\begin{proposition}[Powers of the canonical transition matrix]
\label{prop:powers-canonical-P}
Let $P = \begin{pmatrix} \II_r & \b0_{rm} \\ R & Q \end{pmatrix}$ be in canonical block form. For every integer $k \ge 1$,
\[
P^k = \begin{pmatrix}
\II_r & \b0_{rm} \\
(\II_m + Q + Q^2 + \cdots + Q^{k-1})R & Q^k
\end{pmatrix} = \begin{pmatrix}
\II_r & \b0_{rm} \\
\left(\sum_{j=0}^{k-1} Q^j\right) R & Q^k
\end{pmatrix}
\]
\end{proposition}
\vfill
The $(i,j)$ entry of $Q^k$ is the probability that the chain is in the $j$th transient state after $k$ steps, given that it started in the $i$th transient state
\end{frame}

\begin{frame}{Proof of Proposition~\ref{prop:powers-canonical-P}}
We proceed by mathematical induction on $k \ge 1$
\vfill
\textbf{Base case ($k=1$):} Since $Q^0 = \II_m$, the formula gives $\begin{pmatrix} \II_r & \b0_{rm} \\ \II_m R & Q^1 \end{pmatrix} = P$, which holds
\vfill
\textbf{Inductive step:} Suppose the formula holds for $k \ge 1$. Using block matrix multiplication for $P^{k+1} = P^k P$:
\begin{align*}
P^{k+1} &= \begin{pmatrix}
\II_r & \b0_{rm} \\
(\II_m + Q + \cdots + Q^{k-1})R & Q^k
\end{pmatrix}
\begin{pmatrix}
\II_r & \b0_{rm} \\
R & Q
\end{pmatrix} \\
&= \begin{pmatrix}
\II_r \II_r + \b0_{rm} R & \II_r \b0_{rm} + \b0_{rm} Q \\
(\II_m + Q + \cdots + Q^{k-1})R\II_r + Q^k R & (\II_m + Q + \cdots + Q^{k-1})R\b0_{rm} + Q^k Q
\end{pmatrix} \\
&= \begin{pmatrix}
\II_r & \b0_{rm} \\
(\II_m + Q + Q^2 + \cdots + Q^k)R & Q^{k+1}
\end{pmatrix}
\end{align*}
This completes the induction \hfill $\blacksquare$
\end{frame}


%%%%%%%%%%%%%%%%%%%%%%%%%%%%%%%%%
\SSsubsection{The Neumann series and fundamental matrix}

\begin{frame}{Spectral radius of the transient block $Q$}
\begin{theorem}[Contraction of the transient block]
\label{th:rho-Q-less-than-1}
In any absorbing Markov chain with canonical transient block $Q \in \M_m$, the spectral radius satisfies
\[
\rho(Q) < 1 \quad \text{and consequently} \quad \lim_{k \to \infty} Q^k = \b0_m
\]
\end{theorem}
\vfill
Intuitively, because every transient state has a directed walk to some absorbing state, probability mass continually leaks out of the transient states as $k$ increases
\end{frame}

\begin{frame}{Why $\rho(Q) < 1$: proof sketch}
\begin{itemize}
\item By definition of an absorbing chain, from every transient state $i \in \{1,\ldots,m\}$ there is a directed walk to some absorbing state of length at most $m$
\vfill
\item Thus, starting from transient state $i$, the probability of being absorbed within $m$ steps is strictly positive ($> 0$), so the probability of remaining in the transient states after $m$ steps is strictly less than $1$
\vfill
\item This means every row sum of $Q^m$ is strictly less than $1$:
\[
\|Q^m\|_\infty = \max_{1 \le i \le m} \sum_{j=1}^m (Q^m)_{ij} < 1
\]
\vfill
\item Therefore $\rho(Q)^m = \rho(Q^m) \le \|Q^m\|_\infty < 1$, which implies $\rho(Q) < 1$ and $\lim_{k\to\infty} Q^k = \b0_m$
\end{itemize}
\end{frame}

\begin{frame}{Absorption occurs with probability 1}
Because $(Q^k)_{ij}$ is the probability of being in transient state $j$ at step $k$ starting from transient state $i$, Theorem~\ref{th:rho-Q-less-than-1} immediately answers our first question
\vfill
\begin{importanttheorem}[Eventual absorption]
In an absorbing Markov chain, the probability that the process eventually reaches an absorbing state is $1$, regardless of the initial state
\end{importanttheorem}
\vfill
Indeed, starting from transient state $i$, the probability of still being in \emph{some} transient state after $k$ steps is the $i$th row sum $\sum_{j=1}^m (Q^k)_{ij} = (Q^k \11_m)_i \to 0$ as $k \to \infty$
\end{frame}

\begin{frame}{The Neumann series for matrices}
Recall the scalar geometric series: for $q \in \IR$ with $|q| < 1$, $(1 - q)^{-1} = \sum_{k=0}^\infty q^k$. The exact same identity holds for square matrices!
\vfill
\begin{theorem}[Neumann series]
\label{th:neumann-series}
Let $Q \in \M_m$ be a matrix with spectral radius $\rho(Q) < 1$ (equivalently, $\lim_{k\to\infty} Q^k = \b0_m$). Then $\II_m - Q$ is invertible and its inverse is given by the convergent series
\[
(\II_m - Q)^{-1} = \sum_{k=0}^\infty Q^k = \II_m + Q + Q^2 + Q^3 + \cdots
\]
\end{theorem}
\vfill
In particular, since $Q \ge 0$ for an absorbing Markov chain, every term $Q^k \ge 0$ is nonnegative, so $(\II_m - Q)^{-1} \ge \II_m \ge 0$ has strictly nonnegative entries (and diagonal entries $\ge 1$)
\end{frame}

\begin{frame}{Proof of Theorem~\ref{th:neumann-series} (Neumann series)}
Consider the partial sum $S_k = \II_m + Q + Q^2 + \cdots + Q^{k-1}$ for $k \ge 1$. Multiplying by $\II_m - Q$ gives a telescoping sum:
\begin{align*}
(\II_m - Q) S_k &= (\II_m - Q)(\II_m + Q + Q^2 + \cdots + Q^{k-1}) \\
&= (\II_m + Q + Q^2 + \cdots + Q^{k-1}) - (Q + Q^2 + \cdots + Q^{k-1} + Q^k) \\
&= \II_m - Q^k
\end{align*}
\vfill
First, suppose $(\II_m - Q)\bx = \b0$. Then $\bx = Q\bx = Q^2\bx = \cdots = Q^k\bx \to \b0$ as $k \to \infty$ since $\rho(Q) < 1$. Hence $\nullspace(\II_m - Q) = \{\b0\}$, so $\II_m - Q$ is invertible
\vfill
Multiplying $(\II_m - Q)S_k = \II_m - Q^k$ on the left by $(\II_m - Q)^{-1}$ gives
\[
S_k = (\II_m - Q)^{-1}(\II_m - Q^k) = (\II_m - Q)^{-1} - (\II_m - Q)^{-1}Q^k
\]
Taking the limit as $k \to \infty$ and using $\lim_{k\to\infty} Q^k = \b0_m$ yields
\[
\sum_{j=0}^\infty Q^j = \lim_{k\to\infty} S_k = (\II_m - Q)^{-1} \qquad \blacksquare
\]
\end{frame}

\begin{frame}{The fundamental matrix}
\begin{importantdefinition}[Fundamental matrix]
For an absorbing Markov chain with canonical transient block $Q \in \M_m$, the matrix
\[
N = (\II_m - Q)^{-1} \in \M_m \qquad (= \II_m + Q + Q^2 + \cdots)
\]
is called the \defword{fundamental matrix} of the absorbing Markov chain
\end{importantdefinition}
\vfill
\begin{itemize}
\item By Theorem~\ref{th:rho-Q-less-than-1} and Theorem~\ref{th:neumann-series}, $\II_m - Q$ is always invertible and $N \ge \II_m \ge 0$
\item As we now show, the three quantitative questions about an absorbing chain (expected visits, expected time to absorption and absorption probabilities) can all be answered directly from $N$!
\end{itemize}
\end{frame}

\begin{frame}{Absorption quantities from the fundamental matrix}
\begin{theorem}[Quantitative analysis of absorbing Markov chains]
\label{th:absorbing-chain-quantities}
Let $P = \begin{pmatrix} \II_r & \b0_{rm} \\ R & Q \end{pmatrix}$ be in canonical form with fundamental matrix $N = [n_{ij}] = (\II_m - Q)^{-1}$. Then:
\begin{enumerate}
\item \textbf{Expected visits to transient states:} $n_{ij}$ is the expected total number of times the chain visits the $j$th transient state, given that it starts in the $i$th transient state
\item \textbf{Expected time to absorption:} Let $\bt = N\11_m = (t_1, \ldots, t_m)^T$. Then $t_i = \sum_{j=1}^m n_{ij}$ (the $i$th row sum of $N$) is the expected number of steps before absorption when starting in the $i$th transient state
\item \textbf{Absorption probabilities:} Let $B = [b_{ij}] = N R \in \M_{mr}$. Then $b_{ij}$ is the probability of eventually being absorbed in the $j$th absorbing state when starting in the $i$th transient state
\end{enumerate}
\end{theorem}
\end{frame}

\begin{frame}{Proof of Theorem~\ref{th:absorbing-chain-quantities} (part 1: visits and time)}
\textbf{1. Expected visits $N = (\II_m - Q)^{-1}$:} Starting in transient state $i$, let $X_{ij}^{(k)} \in \{0,1\}$ be the indicator random variable that equals $1$ if the chain is in transient state $j$ at step $k \ge 0$, and $0$ otherwise. Then
\[
\IE[X_{ij}^{(k)}] = \IP(\text{in transient state } j \text{ at step } k \mid \text{start in } i) = (Q^k)_{ij}
\]
By linearity of expectation, the expected total number of visits to transient state $j$ is
\[
\IE\left[\sum_{k=0}^\infty X_{ij}^{(k)}\right] = \sum_{k=0}^\infty (Q^k)_{ij} = \left(\sum_{k=0}^\infty Q^k\right)_{ij} = n_{ij}
\]
\vfill
\textbf{2. Time to absorption $\bt = N\11_m$:} The total number of steps spent in \emph{all} transient states before absorption starting from $i$ is the sum over all $j=1,\ldots,m$ of visits to transient state $j$:
\[
t_i = \sum_{j=1}^m n_{ij} = (N\11_m)_i
\]
\end{frame}

\begin{frame}{Proof of Theorem~\ref{th:absorbing-chain-quantities} (part 2: absorption probabilities)}
\textbf{3. Absorption probabilities $B = N R$:} By Proposition~\ref{prop:powers-canonical-P}, for every $k \ge 1$,
\[
P^k = \begin{pmatrix}
\II_r & \b0_{rm} \\
\left(\sum_{j=0}^{k-1} Q^j\right) R & Q^k
\end{pmatrix}
\]
\vfill
Taking the limit as $k \to \infty$ and using $\sum_{j=0}^\infty Q^j = (\II_m - Q)^{-1} = N$ and $\lim_{k\to\infty} Q^k = \b0_m$ gives
\[
\lim_{k\to\infty} P^k = \begin{pmatrix} \II_r & \b0_{rm} \\ \left(\sum_{j=0}^\infty Q^j\right) R & \lim_{k\to\infty} Q^k \end{pmatrix} = \begin{pmatrix} \II_r & \b0_{rm} \\ N R & \b0_m \end{pmatrix} = \begin{pmatrix} \II_r & \b0_{rm} \\ B & \b0_m \end{pmatrix} \qquad \blacksquare
\]
\end{frame}


%%%%%%%%%%%%%%%%%%%%%%%%%%%%%%%%%
\SSsubsection{Examples and R implementation}

\begin{frame}{Example 1: back to the genetic example}
In the modified orange cat breeding chain ($S_1 = X^O X^O$, $S_2 = X^o X^o$, $S_3 = X^O X^o$), state $S_1$ is absorbing ($r=1$) and $S_2, S_3$ are transient ($m=2$). The matrix is already in canonical form:
\[
P = \begin{pmatrix}
1 & 0 & 0 \\
0 & 0 & 1 \\
1/2 & 0 & 1/2
\end{pmatrix}
= \begin{pmatrix}
\II_1 & \b0_{1,2} \\
R & Q
\end{pmatrix}
\quad \text{with} \quad
R = \begin{pmatrix} 0 \\ 1/2 \end{pmatrix}, \quad
Q = \begin{pmatrix} 0 & 1 \\ 0 & 1/2 \end{pmatrix}
\]
\vfill
We compute $\II_2 - Q$ and invert it using $\begin{pmatrix} a & b \\ c & d \end{pmatrix}^{-1} = \frac{1}{ad-bc}\begin{pmatrix} d & -b \\ -c & a \end{pmatrix}$:
\[
\II_2 - Q = \begin{pmatrix} 1 & -1 \\ 0 & 1/2 \end{pmatrix}
\implies
N = (\II_2 - Q)^{-1} = \frac{1}{1/2} \begin{pmatrix} 1/2 & 1 \\ 0 & 1 \end{pmatrix} = \begin{pmatrix} 1 & 2 \\ 0 & 2 \end{pmatrix}
\]
\end{frame}

\begin{frame}{Interpreting Example 1}
The fundamental matrix for transient states $(S_2, S_3) = (X^o X^o, X^O X^o)$ is
\[
N = \begin{pmatrix} 1 & 2 \\ 0 & 2 \end{pmatrix}
\]
\begin{itemize}
\item Starting with a black female ($S_2$, row 1 of $N$): expected visits are $n_{11} = 1$ generation in $S_2$ (at $t=0$) and $n_{12} = 2$ generations in $S_3$ (tortoiseshell)
\item Starting with a tortoiseshell female ($S_3$, row 2 of $N$): $n_{21} = 0$ generations in $S_2$ and $n_{22} = 2$ generations in $S_3$
\end{itemize}
\vfill
Expected number of generations before absorption into $S_1$ ($X^O X^O$):
\[
\bt = N\11_2 = \begin{pmatrix} 1 & 2 \\ 0 & 2 \end{pmatrix} \begin{pmatrix} 1 \\ 1 \end{pmatrix} = \begin{pmatrix} 3 \\ 2 \end{pmatrix}
\]
\vfill
Absorption probability matrix (into the sole absorbing state $S_1$):
\[
B = N R = \begin{pmatrix} 1 & 2 \\ 0 & 2 \end{pmatrix} \begin{pmatrix} 0 \\ 1/2 \end{pmatrix} = \begin{pmatrix} 1 \\ 1 \end{pmatrix}
\]
\end{frame}

\begin{frame}{Example 2: 4-state chain with 2 absorbing states}
Consider a 4-state absorbing chain where $S_1$ (pass/graduated) and $S_2$ (withdraw) are absorbing ($r=2$), and $S_3$ (Year 1) and $S_4$ (Year 2) are transient ($m=2$):
\vfill
\begin{columns}[c]
\begin{column}{0.48\textwidth}
\[
P = \begin{array}{c@{\hspace{4pt}}c}
& \begin{array}{cccc} S_1 & S_2 & S_3 & S_4 \end{array} \\
\begin{array}{c} S_1 \\ S_2 \\ S_3 \\ S_4 \end{array} &
\left(\begin{array}{cc|cc}
1 & 0 & 0 & 0 \\
0 & 1 & 0 & 0 \\ \hline
0 & 1/4 & 1/4 & 1/2 \\
1/2 & 1/4 & 0 & 1/4
\end{array}\right)
\end{array}
\]
\end{column}
\begin{column}{0.52\textwidth}
\centering
\begin{tikzpicture}[
    ->,
    >=Latex,
    node distance=1.8cm,
    every path/.style={thick},
    abs/.style={state, fill=red!20, minimum size=0.85cm},
    trans/.style={state, fill=black!10, minimum size=0.85cm}
]
    \node[abs] (S2) at (0, 1.8) {$S_2$};
    \node[abs] (S1) at (3.2, 1.8) {$S_1$};
    \node[trans] (S3) at (0, 0) {$S_3$};
    \node[trans] (S4) at (3.2, 0) {$S_4$};
    \path (S1) edge [loop right] node[right] {$1$} (S1);
    \path (S2) edge [loop left] node[left] {$1$} (S2);
    \path (S3) edge [loop left] node[left] {$1/4$} (S3);
    \path (S4) edge [loop right] node[right] {$1/4$} (S4);
    \path (S3) edge node[left] {$1/4$} (S2);
    \path (S3) edge node[below] {$1/2$} (S4);
    \path (S4) edge node[right] {$1/2$} (S1);
    \path (S4) edge node[above,sloped] {$1/4$} (S2);
\end{tikzpicture}
\end{column}
\end{columns}
\vfill
Here the canonical blocks are
\[
R = \begin{pmatrix} 0 & 1/4 \\ 1/2 & 1/4 \end{pmatrix} \in \M_2, \qquad
Q = \begin{pmatrix} 1/4 & 1/2 \\ 0 & 1/4 \end{pmatrix} \in \M_2
\]
\end{frame}

\begin{frame}{Example 2: computing the fundamental matrix $N$}
\textbf{Step 1 (Fundamental matrix $N$):} Compute $\II_2 - Q$ and its determinant:
\[
\II_2 - Q = \begin{pmatrix} 3/4 & -1/2 \\ 0 & 3/4 \end{pmatrix}, \qquad
\det(\II_2 - Q) = \frac{3}{4}\cdot\frac{3}{4} - 0 = \frac{9}{16}
\]
\vfill
Inverting via the $2 \times 2$ inverse formula gives
\[
N = (\II_2 - Q)^{-1} = \frac{16}{9} \begin{pmatrix} 3/4 & 1/2 \\ 0 & 3/4 \end{pmatrix} = \begin{pmatrix} 4/3 & 8/9 \\ 0 & 4/3 \end{pmatrix}
\]
\end{frame}

\begin{frame}{Example 2: expected absorption time $\bt$ and probabilities $B$}
\textbf{Step 2 (Expected steps to absorption $\bt = N\11_2$):}
\[
\bt = \begin{pmatrix} 4/3 & 8/9 \\ 0 & 4/3 \end{pmatrix} \begin{pmatrix} 1 \\ 1 \end{pmatrix} = \begin{pmatrix} 20/9 \\ 4/3 \end{pmatrix} \quad \text{(starting from } S_3 \text{ and } S_4\text{)}
\]
\vfill
\textbf{Step 3 (Absorption probabilities $B = N R$):}
\[
B = N R = \begin{pmatrix} 4/3 & 8/9 \\ 0 & 4/3 \end{pmatrix} \begin{pmatrix} 0 & 1/4 \\ 1/2 & 1/4 \end{pmatrix} = \begin{pmatrix} 4/9 & 5/9 \\ 2/3 & 1/3 \end{pmatrix}
\]
\vfill
Notice each row of $B$ sums to $1$: starting in $S_3$, a student graduates ($S_1$) with probability $4/9$ and withdraws ($S_2$) with probability $5/9$
\end{frame}

\begin{frame}{Example 3: symmetric 5-state random walk}
Return to Motivating example 1 on $\{S_1, S_2, S_3, S_4, S_5\}$ with $p = q = 1/2$ and absorbing walls $S_1, S_5$. Reordering states as $(S_1, S_5, S_2, S_3, S_4)$ puts $P$ in canonical form with $r=2$ and $m=3$:
\[
R = \begin{array}{c@{\hspace{4pt}}c}
& \begin{array}{cc} S_1 & S_5 \end{array} \\
\begin{array}{c} S_2 \\ S_3 \\ S_4 \end{array} &
\left(\begin{array}{cc}
1/2 & 0 \\
0 & 0 \\
0 & 1/2
\end{array}\right)
\end{array}, \qquad
Q = \begin{array}{c@{\hspace{4pt}}c}
& \begin{array}{ccc} S_2 & S_3 & S_4 \end{array} \\
\begin{array}{c} S_2 \\ S_3 \\ S_4 \end{array} &
\left(\begin{array}{ccc}
0 & 1/2 & 0 \\
1/2 & 0 & 1/2 \\
0 & 1/2 & 0
\end{array}\right)
\end{array}
\]
\vfill
Inverting $\II_3 - Q = \begin{pmatrix} 1 & -1/2 & 0 \\ -1/2 & 1 & -1/2 \\ 0 & -1/2 & 1 \end{pmatrix}$ gives
\[
N = (\II_3 - Q)^{-1} = \begin{pmatrix}
3/2 & 1 & 1/2 \\
1 & 2 & 1 \\
1/2 & 1 & 3/2
\end{pmatrix}, \quad
\bt = N\11_3 = \begin{pmatrix} 3 \\ 4 \\ 3 \end{pmatrix}, \quad
B = N R = \begin{pmatrix}
3/4 & 1/4 \\
1/2 & 1/2 \\
1/4 & 3/4
\end{pmatrix}
\]
\end{frame}

\begin{frame}[fragile]{R implementation: computing the fundamental matrix}
We can verify Example 3 directly in base R using \code{solve}:
\begin{knitrout}
\definecolor{shadecolor}{rgb}{0.969, 0.969, 0.969}\color{fgcolor}\begin{kframe}
\begin{alltt}
\hldef{R_mat} \hlkwb{<-} \hlkwd{matrix}\hldef{(}\hlkwd{c}\hldef{(}\hlnum{0.5}\hldef{,} \hlnum{0}\hldef{,}  \hlnum{0}\hldef{,} \hlnum{0}\hldef{,}  \hlnum{0}\hldef{,} \hlnum{0.5}\hldef{),} \hlkwc{nrow} \hldef{=} \hlnum{3}\hldef{,} \hlkwc{byrow} \hldef{=} \hlnum{TRUE}\hldef{,}
                \hlkwc{dimnames} \hldef{=} \hlkwd{list}\hldef{(}\hlkwd{c}\hldef{(}\hlsng{"S2"}\hldef{,} \hlsng{"S3"}\hldef{,} \hlsng{"S4"}\hldef{),} \hlkwd{c}\hldef{(}\hlsng{"S1"}\hldef{,} \hlsng{"S5"}\hldef{)))}
\hldef{Q_mat} \hlkwb{<-} \hlkwd{matrix}\hldef{(}\hlkwd{c}\hldef{(}\hlnum{0}\hldef{,} \hlnum{0.5}\hldef{,} \hlnum{0}\hldef{,}  \hlnum{0.5}\hldef{,} \hlnum{0}\hldef{,} \hlnum{0.5}\hldef{,}  \hlnum{0}\hldef{,} \hlnum{0.5}\hldef{,} \hlnum{0}\hldef{),} \hlkwc{nrow} \hldef{=} \hlnum{3}\hldef{,}
                \hlkwc{dimnames} \hldef{=} \hlkwd{list}\hldef{(}\hlkwd{c}\hldef{(}\hlsng{"S2"}\hldef{,} \hlsng{"S3"}\hldef{,} \hlsng{"S4"}\hldef{),} \hlkwd{c}\hldef{(}\hlsng{"S2"}\hldef{,} \hlsng{"S3"}\hldef{,} \hlsng{"S4"}\hldef{)))}
\hldef{N_mat} \hlkwb{<-} \hlkwd{solve}\hldef{(}\hlkwd{diag}\hldef{(}\hlnum{3}\hldef{)} \hlopt{-} \hldef{Q_mat)}
\hldef{N_mat}
\end{alltt}
\begin{verbatim}
##     S2 S3  S4
## S2 1.5  1 0.5
## S3 1.0  2 1.0
## S4 0.5  1 1.5
\end{verbatim}
\end{kframe}
\end{knitrout}
\end{frame}

\begin{frame}[fragile]{R implementation: absorption times and probabilities}
From \code{N\_mat}, we obtain the expected absorption times $\bt = N\11_3$ and absorption probabilities $B = N R$:
\begin{knitrout}
\definecolor{shadecolor}{rgb}{0.969, 0.969, 0.969}\color{fgcolor}\begin{kframe}
\begin{alltt}
\hldef{t_vec} \hlkwb{<-} \hlkwd{rowSums}\hldef{(N_mat)}
\hldef{t_vec}
\end{alltt}
\begin{verbatim}
## S2 S3 S4 
##  3  4  3
\end{verbatim}
\begin{alltt}
\hldef{B_mat} \hlkwb{<-} \hldef{N_mat} \hlopt{%*%} \hldef{R_mat}
\hldef{B_mat}
\end{alltt}
\begin{verbatim}
##      S1   S5
## S2 0.75 0.25
## S3 0.50 0.50
## S4 0.25 0.75
\end{verbatim}
\end{kframe}
\end{knitrout}
\end{frame}

\begin{frame}[fragile]{Using the \code{markovchain} package}
The \code{markovchain} package uses the same row-stochastic convention, so we pass \code{P\_abs} directly:
\begin{knitrout}
\definecolor{shadecolor}{rgb}{0.969, 0.969, 0.969}\color{fgcolor}\begin{kframe}
\begin{alltt}
\hlkwd{library}\hldef{(markovchain)}
\hldef{P_abs} \hlkwb{<-} \hlkwd{matrix}\hldef{(}\hlkwd{c}\hldef{(}\hlnum{1}\hldef{,}   \hlnum{0}\hldef{,} \hlnum{0}\hldef{,}
                  \hlnum{0}\hldef{,}   \hlnum{0}\hldef{,} \hlnum{1}\hldef{,}
                  \hlnum{0.5}\hldef{,} \hlnum{0}\hldef{,} \hlnum{0.5}\hldef{),} \hlkwc{nrow} \hldef{=} \hlnum{3}\hldef{,} \hlkwc{byrow} \hldef{=} \hlnum{TRUE}\hldef{)}
\hldef{mc_abs} \hlkwb{<-} \hlkwd{new}\hldef{(}\hlsng{"markovchain"}\hldef{,}
              \hlkwc{states} \hldef{=} \hlkwd{sprintf}\hldef{(}\hlsng{"S_%d"}\hldef{,} \hlnum{1}\hlopt{:}\hlnum{3}\hldef{),}
              \hlkwc{transitionMatrix} \hldef{= P_abs,}
              \hlkwc{name} \hldef{=} \hlsng{"orange_abs"}\hldef{)}
\end{alltt}
\end{kframe}
\end{knitrout}
\end{frame}

\begin{frame}[fragile]{Inspecting the chain with \code{summary}}
\begin{knitrout}
\definecolor{shadecolor}{rgb}{0.969, 0.969, 0.969}\color{fgcolor}\begin{kframe}
\begin{alltt}
\hlkwd{summary}\hldef{(mc_abs)}
\end{alltt}
\begin{verbatim}
## orange_abs  Markov chain that is composed by: 
## Closed classes: 
## S_1 
## Recurrent classes: 
## {S_1}
## Transient classes: 
## {S_2},{S_3}
## The Markov chain is not irreducible 
## The absorbing states are: S_1
\end{verbatim}
\end{kframe}
\end{knitrout}
\end{frame}

\begin{frame}[fragile]{Absorption times and probabilities in \code{markovchain}}
\begin{knitrout}
\definecolor{shadecolor}{rgb}{0.969, 0.969, 0.969}\color{fgcolor}\begin{kframe}
\begin{alltt}
\hlkwd{absorbingStates}\hldef{(mc_abs)}
\end{alltt}
\begin{verbatim}
## [1] "S_1"
\end{verbatim}
\begin{alltt}
\hlkwd{meanAbsorptionTime}\hldef{(mc_abs)}
\end{alltt}
\begin{verbatim}
## S_2 S_3 
##   3   2
\end{verbatim}
\begin{alltt}
\hlkwd{absorptionProbabilities}\hldef{(mc_abs)}
\end{alltt}
\begin{verbatim}
##     S_1
## S_2   1
## S_3   1
\end{verbatim}
\end{kframe}
\end{knitrout}
\end{frame}

\begin{frame}[fragile]{Simulating a realisation with \code{DTMCPack}}
\begin{knitrout}
\definecolor{shadecolor}{rgb}{0.969, 0.969, 0.969}\color{fgcolor}\begin{kframe}
\begin{alltt}
\hlkwd{library}\hldef{(DTMCPack)}
\hlkwd{set.seed}\hldef{(}\hlnum{2740}\hldef{)}
\hldef{IC} \hlkwb{<-} \hlkwd{c}\hldef{(}\hlnum{0}\hldef{,} \hlnum{1}\hldef{,} \hlnum{0}\hldef{)} \hlcom{# Start in S_2 (black female)}
\hldef{sol} \hlkwb{<-} \hlkwd{DTMC}\hldef{(P_abs, IC,} \hlnum{10}\hldef{,} \hlkwc{trace} \hldef{=} \hlnum{TRUE}\hldef{)}
\end{alltt}
\begin{verbatim}
## NULL
\end{verbatim}
\end{kframe}
\end{knitrout}
\end{frame}

\maxFrameImage{FIGS/U09-U09-absorbing-realisation-1.pdf}

\begin{frame}{Example 4: ``limited'' gambler's ruin (\$500K start, \$1M cap)}
Suppose a gambler starts with $\$500\,000$ and places $\$1$ fair bets ($p = q = 1/2$), stopping either at ruin ($\$0$) or when the bank halts betting at $M = \$1\,000\,000$
\begin{itemize}
\item Absorbing states: $0$ and $M = 10^6$; transient states: $\{1, 2, \ldots, 999\,999\}$ ($m = 999\,999$)
\item Forming $N = (\II_m - Q)^{-1}$ as a dense matrix is impossible: $m^2 \approx 10^{12}$ entries ($\approx 8\text{ TB}$ of RAM)
\end{itemize}
\vfill
Instead, $\bt = (\II_m - Q)^{-1}\11_m$ solves the \textbf{tridiagonal linear system} $(\II_m - Q)\bt = \11_m$, where row $i$ reads
\[
t_i = 1 + \frac{1}{2} t_{i-1} + \frac{1}{2} t_{i+1}, \qquad t_0 = t_M = 0
\]
This recurrence has the closed-form solution $t_k = k(M - k)$, so starting at $k = 500\,000$:
\[
t_{500\,000} = (5 \times 10^5)(10^6 - 5 \times 10^5) = 2.5 \times 10^{11} \text{ steps}
\]
By symmetry ($k = M/2$, $p = q = 1/2$), ruin occurs with probability $1/2$ and the expected number of steps until ruin (given ruin occurs) is also $2.5 \times 10^{11}$
\end{frame}

\begin{frame}[fragile]{Sparse R implementation for the limited gambler's ruin}
Using \code{Matrix::bandSparse}, we store only the three nonzero diagonals of $\II_m - Q$ ($O(m)$ memory) and solve $(\II_m - Q)\bt = \11_m$ directly:
\begin{knitrout}
\definecolor{shadecolor}{rgb}{0.969, 0.969, 0.969}\color{fgcolor}\begin{kframe}
\begin{alltt}
\hlkwd{library}\hldef{(Matrix)}
\hldef{M} \hlkwb{<-} \hlnum{1e6}
\hldef{m} \hlkwb{<-} \hldef{M} \hlopt{-} \hlnum{1}
\hldef{p} \hlkwb{<-} \hlnum{0.5}
\hldef{q} \hlkwb{<-} \hlnum{1} \hlopt{-} \hldef{p}
\hldef{A} \hlkwb{<-} \hlkwd{bandSparse}\hldef{(m,} \hlkwc{k} \hldef{=} \hlkwd{c}\hldef{(}\hlopt{-}\hlnum{1}\hldef{,} \hlnum{0}\hldef{,} \hlnum{1}\hldef{),}
                \hlkwc{diagonals} \hldef{=} \hlkwd{list}\hldef{(}\hlkwd{rep}\hldef{(}\hlopt{-}\hldef{q, m} \hlopt{-} \hlnum{1}\hldef{),}
                                 \hlkwd{rep}\hldef{(}\hlnum{1}\hldef{, m),}
                                 \hlkwd{rep}\hldef{(}\hlopt{-}\hldef{p, m} \hlopt{-} \hlnum{1}\hldef{)))}
\hldef{t_vec} \hlkwb{<-} \hlkwd{solve}\hldef{(A,} \hlkwd{rep}\hldef{(}\hlnum{1}\hldef{, m))}
\hldef{t_vec[}\hlnum{5e5}\hldef{]}
\end{alltt}
\begin{verbatim}
## [1] 250000149671
\end{verbatim}
\end{kframe}
\end{knitrout}
\end{frame}


%%%%%%%%%%%%%%%%%%%%%%%%%%%%%%%%%
%%%%%%%%%%%%%%%%%%%%%%%%%%%%%%%%%
%% SECTION 2: SPECTRAL CENTRALITY IN NETWORKS
%%%%%%%%%%%%%%%%%%%%%%%%%%%%%%%%%
%%%%%%%%%%%%%%%%%%%%%%%%%%%%%%%%%
\SSsection{Spectral centrality in networks}

\SSsubsection{Degree and eigenvector centrality}

\begin{frame}{Why spectral centrality?}
Combinatorial centrality measures on graphs are often based on counting shortest paths (such as closeness and betweenness centrality)
\vfill
\textbf{Computational bottleneck on massive networks:}
\begin{itemize}
\item All-pairs shortest-path algorithms require $O(n m)$ or $O(n^3)$ operations
\item For a social network or web graph with $n \sim 10^8$ vertices, $n^2 \sim 10^{16}$ is completely intractable
\end{itemize}
\vfill
\textbf{The spectral alternative:}
\begin{itemize}
\item Express vertex importance via matrix-vector products $A^T\bx$, eigenvectors $A^T\bx = \lambda\bx$ or linear systems $(\II - \alpha A^T)\bx = \bb$
\item When the adjacency matrix $A$ is \textbf{sparse} ($m = |E| \ll n^2$ non-zero entries), one sparse matrix-vector product $A^T\bx$ takes only $O(m)$ time!
\end{itemize}
\end{frame}

\begin{frame}{Convention for directed adjacency matrices}
Using our standard row-to-column convention from Unit 08, the adjacency matrix $A = [a_{ij}] \in \M_n$ of a directed graph $\G = (V, E)$ on $n$ vertices has entries:
\vfill
\[
a_{ij} = \begin{cases}
1 & \text{if there is a directed arc } i \to j \text{ (from } i \text{ to } j\text{)} \\
0 & \text{otherwise}
\end{cases}
\]
\vfill
With this convention:
\begin{itemize}
\item \textbf{Row $i$} of $A$ records arcs \emph{leaving} vertex $i$, so the \defword{outdegree} of vertex $i$ is $d_i^+ = \sum_{j=1}^n a_{ij} = (A\11_n)_i$
\item \textbf{Column $j$} of $A$ records arcs \emph{entering} vertex $j$, so the \defword{indegree} of vertex $j$ is $d_j^- = \sum_{i=1}^n a_{ij} = (A^T\11_n)_j = (\11_n^T A)_j$
\item For an undirected graph, $A = A^T$ is symmetric and $d_i^- = d_i^+ = d_i$
\end{itemize}
\end{frame}

\begin{frame}{First-order measure: indegree centrality}
The simplest measure of prestige in a directed network is \defword{indegree centrality}:
\[
\bd^- = A^T \11_n, \qquad d_i^- = \sum_{j=1}^n a_{ji}
\]
Each incoming arc $j \to i$ is treated as one ``vote'' for the importance of vertex $i$
\vfill
\textbf{Fundamental limitation:}
\begin{itemize}
\item Indegree treats all incoming links as equal
\item Receiving a link from a major hub or authoritative page counts the exact same ($+1$) as receiving a link from an isolated spam page
\item We want a \textbf{recursive} notion of prestige: \emph{a vertex is important if it is pointed to by other important vertices}
\end{itemize}
\end{frame}

\begin{frame}{Eigenvector centrality}
Instead of summing $1$ for each incoming neighbour $j \to i$, we sum the \emph{importance scores} $x_j$ of those in-neighbours and require $x_i$ to be proportional to that sum
\vfill
\begin{importantdefinition}[Eigenvector centrality]
Let $A = [a_{ij}] \in \M_n$ be the adjacency matrix of a graph or digraph (with $a_{ij} = 1$ if $i \to j$). The \defword{eigenvector centrality} (or \defword{prestige score}) of vertex $i$ is the $i$th entry of the nonnegative eigenvector $\bx \ge \b0$ satisfying
\[
\lambda x_i = \sum_{j=1}^n a_{ji} x_j, \quad i=1,\ldots,n \iff A^T \bx = \lambda \bx \quad \bigl(\text{or } \bx^T A = \lambda \bx^T\bigr)
\]
associated with the dominant eigenvalue $\lambda = \rho(A^T) = \rho(A)$
\end{importantdefinition}
\vfill
Notice that $x_i = \frac{1}{\lambda}\sum_{j \in \N^{\text{in}}(i)} x_j$ depends on both the \textbf{quantity} of in-neighbours $|\N^{\text{in}}(i)|$ and the \textbf{quality} $x_j$ of each in-neighbour
\end{frame}

\begin{frame}{Why Perron-Frobenius theory is essential here}
\begin{theorem}[Existence and uniqueness of eigenvector centrality]
Let $\G = (V, E)$ be a strongly connected digraph (or connected undirected graph) with adjacency matrix $A \in \M_n$. Then $A$ (and hence $A^T$) is irreducible and nonnegative, so by the Perron-Frobenius theorem (Unit 08):
\begin{enumerate}
\item The spectral radius $\lambda_1 = \rho(A) > 0$ is a simple eigenvalue of $A$ and $A^T$
\item There exists a unique eigenvector $\bx > \b0$ (the \defword{Perron eigenvector} of $A^T$, or left Perron eigenvector of $A$) satisfying
\[
A^T \bx = \rho(A)\bx \quad \bigl(\text{equivalently } \bx^T A = \rho(A)\bx^T\bigr) \quad \text{and} \quad \11_n^T \bx = \sum_{i=1}^n x_i = 1
\]
\item No other eigenvalue of $A^T$ has a nonnegative eigenvector
\end{enumerate}
\end{theorem}
\end{frame}

\begin{frame}{Constructing a stochastic matrix from an adjacency matrix}
Let $A = [a_{ij}] \in \M_n$ be the adjacency matrix of a simple undirected graph $\G = (V, E)$ with no isolated vertices, and let $D = \diag(d_{11}, \ldots, d_{nn})$ be its degree matrix, with diagonal entries
\[
d_{ii} = \sum_{j=1}^n a_{ij} > 0
\]
\vfill
\begin{proposition}[Random walk transition matrix on an undirected graph]
\label{prop:DinvA-row-stochastic}
The matrix $P = D^{-1} A$ is \textbf{row-stochastic}, and defines the transition matrix of a random walk on $\G$ where a walker at vertex $i$ moves to a uniformly chosen neighbour $j \in \N(i)$ with probability $1/d_{ii}$
\end{proposition}
\end{frame}

\begin{frame}{Proof of Proposition~\ref{prop:DinvA-row-stochastic}}
Write the entries of the diagonal matrix $D^{-1}$ as $d_{ij}^{-1}$, where $d_{ii}^{-1} = 1/d_{ii}$ for $1 \le i \le n$ and $d_{ij}^{-1} = 0$ for $i \neq j$
\vfill
By matrix multiplication, the $(i, j)$ entry of $D^{-1} A$ is
\[
(D^{-1} A)_{ij} = \sum_{k=1}^n d_{ik}^{-1} a_{kj} = d_{ii}^{-1} a_{ij} = \frac{a_{ij}}{d_{ii}} \ge 0, \qquad i, j = 1, \ldots, n
\]
\vfill
Therefore, the sum of row $i$ of $D^{-1} A$ is
\begin{align*}
\sum_{j=1}^n d_{ii}^{-1} a_{ij}
&= d_{ii}^{-1} \sum_{j=1}^n a_{ij} \\
&= d_{ii}^{-1} d_{ii} = 1
\end{align*}
Hence every row sums to $1$ and $D^{-1} A$ is row-stochastic \hfill $\blacksquare$
\end{frame}

\begin{frame}{Computing eigenvector centrality: the power method}
For large networks, computing the full characteristic polynomial is impossible. Instead, we compute the dominant eigenpair $(\lambda_1, \bx)$ using the \defword{power method}:
\vfill
Start with an initial positive vector $\bx^{(0)} > \b0$ (such as $\bx^{(0)} = \frac{1}{n}\11_n$). For $k = 1, 2, \ldots$:
\begin{enumerate}
\item Multiply by $A^T$: $\by^{(k)} = A^T \bx^{(k-1)}$
\item Normalise (using the $\ell_1$ norm $\|\by^{(k)}\|_1 = \11_n^T \by^{(k)}$ or the maximum-magnitude entry $m(\by^{(k)})$):
\[
\bx^{(k)} = \frac{\by^{(k)}}{\|\by^{(k)}\|_1}
\]
\end{enumerate}
until $\|\bx^{(k)} - \bx^{(k-1)}\|_1 < \varepsilon$
\end{frame}

\begin{frame}{Convergence of the power method}
\begin{theorem}[Convergence of the power method]
If $A$ is primitive (so $\lambda_1 = \rho(A) > |\lambda_2| \ge \cdots \ge |\lambda_n|$), then $\bx^{(k)} \to \bx$ as $k \to \infty$, and the error decays at the geometric rate
\[
\|\bx^{(k)} - \bx\| = O\!\left(\left|\frac{\lambda_2}{\lambda_1}\right|^k\right)
\]
\end{theorem}
\vfill
\begin{itemize}
\item Each iteration only requires a sparse matrix-vector product $\by^{(k)} = A^T \bx^{(k-1)}$ ($O(m)$ operations for $m$ edges)
\item The spectral gap $\lambda_1 - |\lambda_2|$ controls how many iterations are needed
\end{itemize}
\end{frame}


%%%%%%%%%%%%%%%%%%%%%%%%%%%%%%%%%
\SSsubsection{Issue with directed graphs and Katz centrality}

\begin{frame}{What goes wrong on general directed graphs?}
Real-world directed networks (citation graphs, web graphs, DAGs) are rarely strongly connected! When $\G(A)$ is \textbf{reducible}, eigenvector centrality suffers from a severe pathology:
\vfill
\begin{itemize}
\item Suppose vertex $i$ is a \textbf{source node} (indegree $d_i^- = 0$, so nobody points to $i$)
\vfill
\item Then column $i$ of $A$ (row $i$ of $A^T$) is all zeros, so the eigenvector equation $A^T\bx = \lambda\bx$ with $\lambda > 0$ gives
\[
\lambda x_i = \sum_{j=1}^n a_{ji} x_j = 0 \implies x_i = 0
\]
\vfill
\item \textbf{Domino effect:} If vertex $k$ only receives links from source nodes with $x_i = 0$, then $\lambda x_k = \sum_{j \to k} x_j = 0 \implies x_k = 0$ as well!
\end{itemize}
\end{frame}

\begin{frame}{Example of the zero-propagation domino effect}
Consider a simple 4-vertex directed chain where only vertex $4$ has a self-loop:
\vfill
\begin{center}
\begin{tikzpicture}[
    ->,
    >=Latex,
    node distance=2cm,
    every path/.style={thick},
    vnode/.style={state, fill=black!10, minimum size=0.95cm}
]
    \node[vnode] (1) {$1$};
    \node[vnode, right=of 1] (2) {$2$};
    \node[vnode, right=of 2] (3) {$3$};
    \node[state, fill=orange!25, minimum size=0.95cm, right=of 3] (4) {$4$};
    \draw[->] (1) to (2);
    \draw[->] (2) to (3);
    \draw[->] (3) to (4);
    \draw[->] (4) to [loop right] (4);
\end{tikzpicture}
\end{center}
\vfill
\begin{itemize}
\item Vertex $1$ has no incoming arcs ($d_1^- = 0$), so $x_1 = 0$
\item Vertex $2$ only receives an arc from vertex $1$, so $\lambda x_2 = x_1 = 0 \implies x_2 = 0$
\item Vertex $3$ only receives an arc from vertex $2$, so $\lambda x_3 = x_2 = 0 \implies x_3 = 0$
\item Only vertex $4$ gets non-zero eigenvector centrality ($x_4 = 1$)!
\end{itemize}
\end{frame}

\begin{frame}{Katz centrality (1953): baseline ``free'' importance}
Leo Katz (1953) resolved the zero-propagation problem on directed graphs by giving every vertex a small baseline ``free'' importance $\beta > 0$, plus damped importance $\alpha > 0$ from its in-neighbours
\vfill
\begin{definition}[Katz centrality]
Let $A \in \M_n$ be the adjacency matrix of a directed graph (with $a_{ij} = 1$ if $i \to j$), $\alpha > 0$ be an \defword{attenuation factor} (or damping parameter) and $\beta > 0$ be a baseline constant. The \defword{Katz centrality} vector $\bx \in \IR^n$ is the solution to
\[
x_i = \alpha \sum_{j=1}^n a_{ji} x_j + \beta, \quad i=1,\ldots,n
\iff
\bx = \alpha A^T \bx + \beta \11_n
\]
\end{definition}
\vfill
Rearranging terms gives the linear system $(\II_n - \alpha A^T)\bx = \beta \11_n$ (or $\bx^T(\II_n - \alpha A) = \beta\11_n^T$). Since $\beta$ only scales $\bx$ by a positive constant, we typically set $\beta = 1$
\end{frame}

\begin{frame}{Existence of Katz centrality via the Neumann series}
When is $\II_n - \alpha A^T$ guaranteed to be invertible with a strictly positive solution $\bx > \b0$?
\vfill
\begin{theorem}[Existence and walk-counting expansion of Katz centrality]
Let $A \in \M_n$ be a nonnegative adjacency matrix and suppose $0 < \alpha < \frac{1}{\rho(A)}$ and $\beta > 0$. Then $\rho(\alpha A^T) = \alpha \rho(A) < 1$, so by Theorem~\ref{th:neumann-series} (the Neumann series), $\II_n - \alpha A^T$ is invertible and
\[
\bx = \beta (\II_n - \alpha A^T)^{-1} \11_n = \beta \left(\sum_{k=0}^\infty \alpha^k (A^T)^k\right) \11_n > \b0
\]
\end{theorem}
\vfill
Expanding the series term by term gives
\[
\bx = \beta \left(\11_n + \alpha A^T \11_n + \alpha^2 (A^T)^2 \11_n + \alpha^3 (A^T)^3 \11_n + \cdots\right)
\]
\end{frame}

\begin{frame}{Walk-counting interpretation of Katz centrality}
What does each term $\alpha^k (A^T)^k \11_n$ in the Neumann series represent?
\vfill
\begin{itemize}
\item Recall from Unit 08 that $(A^k)_{ji}$ is the number of directed walks of length $k$ from vertex $j$ to vertex $i$
\vfill
\item Therefore $((A^T)^k \11_n)_i = \sum_{j=1}^n (A^k)_{ji}$ is the total number of directed walks of length $k$ from \emph{any} vertex $j$ terminating at vertex $i$
\vfill
\item Katz centrality sums walks of \textbf{all lengths} $k \ge 0$ ending at $i$, discounted by $\alpha^k$:
\[
x_i = \beta \big(\underbrace{1}_{k=0} + \alpha \underbrace{d_i^-}_{k=1} + \alpha^2 \underbrace{(\text{\# walks of length } 2 \text{ into } i)}_{k=2} + \cdots\big)
\]
\end{itemize}
\end{frame}

\begin{frame}{Example: Katz centrality on a 3-vertex digraph (setup)}
Consider the 3-vertex directed graph with arcs $1 \to 2$, $1 \to 3$ and $2 \to 3$:
\vfill
\begin{columns}[c]
\begin{column}{0.45\textwidth}
\centering
\begin{tikzpicture}[
    ->,
    >=Latex,
    node distance=2cm,
    every path/.style={thick},
    vnode/.style={state, fill=black!10, minimum size=0.95cm}
]
    \node[vnode] (1) at (0, 1.6) {$1$};
    \node[vnode] (2) at (2.4, 1.6) {$2$};
    \node[vnode] (3) at (1.2, 0) {$3$};
    \draw[->] (1) to (2);
    \draw[->] (1) to (3);
    \draw[->] (2) to (3);
\end{tikzpicture}
\end{column}
\begin{column}{0.55\textwidth}
With $a_{ij} = 1$ if $i \to j$, the adjacency matrix is strictly upper-triangular:
\[
A = \begin{pmatrix}
0 & 1 & 1 \\
0 & 0 & 1 \\
0 & 0 & 0
\end{pmatrix},
\qquad
A^T = \begin{pmatrix}
0 & 0 & 0 \\
1 & 0 & 0 \\
1 & 1 & 0
\end{pmatrix}
\]
Here $\rho(A) = 0$ (the graph is a DAG), so any $\alpha > 0$ satisfies $\rho(\alpha A^T) = 0 < 1$
\end{column}
\end{columns}
\end{frame}

\begin{frame}{Example: solving for Katz centrality on the 3-vertex digraph}
Let us compute Katz centrality with $\alpha = 1/2$ and $\beta = 1$ by solving $(\II_3 - \frac{1}{2}A^T)\bx = \11_3$:
\vfill
\[
\II_3 - \frac{1}{2}A^T = \begin{pmatrix}
1 & 0 & 0 \\
-1/2 & 1 & 0 \\
-1/2 & -1/2 & 1
\end{pmatrix}
\implies
\begin{cases}
x_1 = 1 \\[3pt]
x_2 = \frac{1}{2}x_1 + 1 = \frac{3}{2} \\[3pt]
x_3 = \frac{1}{2}x_1 + \frac{1}{2}x_2 + 1 = \frac{1}{2} + \frac{3}{4} + 1 = \frac{9}{4}
\end{cases}
\]
\vfill
Normalising $\bx = (1, 3/2, 9/4)^T = \frac{1}{4}(4, 6, 9)^T$ by $\11_3^T \bx = 19/4$ gives
\[
\tilde{\bx} = \begin{pmatrix} 4/19 \\ 6/19 \\ 9/19 \end{pmatrix} \implies \text{Ranking: } 3 \succ 2 \succ 1
\]
\end{frame}


%%%%%%%%%%%%%%%%%%%%%%%%%%%%%%%%%
%%%%%%%%%%%%%%%%%%%%%%%%%%%%%%%%%
%% SECTION 3: PAGERANK AND THE GOOGLE MATRIX
%%%%%%%%%%%%%%%%%%%%%%%%%%%%%%%%%
%%%%%%%%%%%%%%%%%%%%%%%%%%%%%%%%%
\SSsection{PageRank and the Google matrix}

\SSsubsection{Motivation and vote splitting}

\begin{frame}{What makes an important webpage?}
In the early days of the World Wide Web (mid-1990s), search engines ranked pages primarily by keyword frequency on the page itself, which was easily manipulated
\vfill
In 1998, Sergey Brin, Larry Page, Rajeev Motwani and Terry Winograd introduced \defword{PageRank}, the algorithm at the heart of Google
\vfill
\textbf{Core idea:} Use the directed hyperlink structure of the Web to measure authority
\begin{itemize}
\item Webpages are vertices $V = \{1, \ldots, n\}$; a hyperlink from page $i$ to page $j$ is a directed arc $i \to j$
\item A link $i \to j$ is a recommendation or ``vote'' from page $i$ for the importance of page $j$
\item \textbf{Random surfer model:} Imagine a web surfer who clicks links at random from page to page. The PageRank of page $i$ is the long-run proportion of time the random surfer spends on page $i$
\end{itemize}
\end{frame}

\begin{frame}{From Katz centrality to vote splitting}
Why not simply use Katz or eigenvector centrality on the Web adjacency matrix $A$?
\vfill
\textbf{The ``generous hub'' flaw of Katz/eigenvector centrality:}
\begin{itemize}
\item In Katz centrality ($x_i = \alpha \sum_{j \to i} x_j + \beta$), if an important directory page $j$ links to $10{,}000$ pages, \emph{every single one} of those $10{,}000$ pages receives $\alpha x_j$ importance!
\item A page with more outgoing links exerts more total influence on the network
\end{itemize}
\vfill
\textbf{Brin and Page's vote-splitting principle:}
\begin{itemize}
\item Each page $j$ has a single unit of voting power equal to its own importance $x_j$
\item If page $j$ has $d_j^+ > 0$ outgoing links, it splits its importance \textbf{equally} among them, passing $x_j / d_j^+$ along each outgoing link:
\[
x_i = \sum_{j \to i} \frac{x_j}{d_j^+}
\]
\end{itemize}
\end{frame}

\begin{frame}{The hyperlink transition matrix}
Let $A = [a_{ij}] \in \M_n$ be the directed adjacency matrix ($a_{ij} = 1$ if page $i$ links to page $j$). Suppose first that every page has at least one outgoing link ($d_i^+ = \sum_{j=1}^n a_{ij} \ge 1$)
\vfill
Define $D^+ = \diag(d_1^+, \ldots, d_n^+)$ and the \defword{hyperlink matrix} $P = [p_{ij}] = (D^+)^{-1} A \in \M_n$:
\[
p_{ij} = \begin{cases}
\dfrac{1}{d_i^+} & \text{if page } i \text{ links to page } j \\
0 & \text{otherwise}
\end{cases}
\]
\vfill
\begin{proposition}
If $d_i^+ \ge 1$ for all $i = 1, \ldots, n$, then $P = (D^+)^{-1} A$ is \textbf{row-stochastic} ($P\11_n = \11_n$), and the vote-splitting equations $x_i = \sum_{j \to i} x_j / d_j^+$ take the fixed-point form
\[
\bx^T P = \bx^T \qquad \bigl(\text{or equivalently } P^T \bx = \bx\bigr)
\]
\end{proposition}
\end{frame}


%%%%%%%%%%%%%%%%%%%%%%%%%%%%%%%%%
\SSsubsection{Obstacles on real web graphs and the Google matrix}

\begin{frame}{Two linear algebra obstacles on real web graphs}
If we try to solve $\bx^T P = \bx^T$ (or $P^T\bx = \bx$) directly on the real Web, two fatal problems occur:
\vfill
\begin{columns}[T]
\begin{column}{0.48\textwidth}
\textbf{Obstacle 1: Dangling nodes ($d_i^+ = 0$)}
\begin{itemize}
\item Many pages (PDF files, images, dead-end pages) have \textbf{no outgoing links}: $d_i^+ = 0$
\item Row $i$ of the raw hyperlink matrix is $\b0^T$, so probability mass ``leaks'' out of the network and $P$ is not stochastic!
\end{itemize}
\begin{center}
\begin{tikzpicture}[->,>=Latex,node distance=1.5cm,every path/.style={thick}]
\node[state,fill=black!10,minimum size=0.8cm] (1) {$1$};
\node[state,fill=black!10,minimum size=0.8cm,right=of 1] (2) {$2$};
\node[state,fill=red!20,minimum size=0.8cm,right=of 2] (3) {$3$};
\draw[->] (1) to [bend left=18] (2);
\draw[->] (2) to [bend left=18] (1);
\draw[->] (2) to (3);
\end{tikzpicture}
\end{center}
\end{column}
\begin{column}{0.48\textwidth}
\textbf{Obstacle 2: Reducibility and spider traps}
\begin{itemize}
\item The Web is not strongly connected ($\G(P)$ is reducible) and contains \textbf{spider traps} (closed strongly connected components) and periodic cycles
\item $P^T\bx = \bx$ may have multiple linearly independent solutions, assign $x_i = 0$ outside traps, or cause power iteration to oscillate!
\end{itemize}
\begin{center}
\begin{tikzpicture}[->,>=Latex,node distance=1.5cm,every path/.style={thick}]
\node[state,fill=black!10,minimum size=0.8cm] (1) {$1$};
\node[state,fill=orange!25,minimum size=0.8cm,right=of 1] (2) {$2$};
\node[state,fill=orange!25,minimum size=0.8cm,right=of 2] (3) {$3$};
\draw[->] (1) to (2);
\draw[->] (2) to [bend left=18] (3);
\draw[->] (3) to [bend left=18] (2);
\end{tikzpicture}
\end{center}
\end{column}
\end{columns}
\end{frame}

\begin{frame}{Fixing Obstacle 1: handling dangling nodes}
Before forming the Google matrix, we must ensure that $P$ is genuinely row-stochastic ($P\11_n = \11_n$)
\vfill
Whenever a vertex $i$ is a \defword{dangling node} ($d_i^+ = 0$, giving a zero row in the raw hyperlink matrix), there are two standard conventions:
\begin{enumerate}
\item \textbf{Uniform teleportation from dead ends:} Replace the zero row $i$ with the uniform probability row vector $\frac{1}{n}\11_n^T = (1/n, \ldots, 1/n)$ (a random surfer reaching a dead-end page picks any page on the Web uniformly at random)
\item \textbf{Self-loop at dead ends:} Add a self-loop $i \to i$ so that $p_{ii} = 1$ (row $i$ becomes $\be_i^T$)
\end{enumerate}
\vfill
In both conventions, every row of the resulting matrix $P$ is a valid probability vector ($P \ge 0$ and $P\11_n = \11_n$). From now on, $P \in \M_n$ denotes this row-stochastic web transition matrix
\end{frame}

\begin{frame}{Fixing Obstacle 2: teleportation and the Google matrix}
To overcome reducibility, spider traps and periodicity, Brin and Page introduced a \defword{damping factor} (or teleportation parameter) $\alpha \in (0, 1)$, typically chosen as $\alpha = 0.85$:
\begin{itemize}
\item With probability $\alpha$, the random surfer follows an outgoing link on the current page (according to $P$)
\item With probability $1 - \alpha$, the surfer gets bored and jumps (teleports) to a page chosen uniformly at random from all $n$ pages (according to $\frac{1}{n}\11_n^T$)
\end{itemize}
\vfill
\begin{importantdefinition}[Google matrix]
Let $P \in \M_n$ be a row-stochastic transition matrix and let $\alpha \in (0, 1)$ be a damping factor. The \defword{Google matrix} is the $n \times n$ matrix
\[
G = \alpha P + (1 - \alpha)\frac{1}{n}\11_n \11_n^T = \alpha P + \frac{1 - \alpha}{n}\IJ_n
\]
where $\IJ_n = \11_n \11_n^T \in \M_n$ is the $n \times n$ matrix of all ones
\end{importantdefinition}
\end{frame}

\begin{frame}{Existence, uniqueness and convergence of PageRank}
Let $P \in \M_n$ be row-stochastic and $\alpha \in (0, 1)$. Then $G = \alpha P + \frac{1-\alpha}{n}\11_n \11_n^T$ is \textbf{row-stochastic} and \textbf{strictly positive} ($G > 0$), hence \textbf{primitive}. By Perron's theorem (Unit 08):
\vfill
\begin{importanttheorem}[PageRank theorem]
\label{th:pagerank-theorem}
\begin{enumerate}
\item The spectral radius $\lambda_1 = \rho(G) = 1$ is a simple, strictly dominant eigenvalue of $G$ (with $|\lambda| \le \alpha < 1$ for every other eigenvalue of $G$)
\item There exists a unique stationary probability vector $\bx > \b0$ (the \defword{PageRank vector}) satisfying
\[
\bx^T G = \bx^T \quad \bigl(\text{equivalently } G^T \bx = \bx\bigr) \quad \text{and} \quad \11_n^T \bx = 1
\]
\item For any initial probability distribution $\bx^{(0)}$, the power iteration $(\bx^{(k)})^T = (\bx^{(0)})^T G^k$ (or $\bx^{(k)} = (G^T)^k \bx^{(0)}$) converges to $\bx$ as $k \to \infty$ at the geometric rate $O(\alpha^k)$
\end{enumerate}
\end{importanttheorem}
\end{frame}

\begin{frame}{Proof of Theorem~\ref{th:pagerank-theorem}}
\textbf{1. Strict positivity ($G > 0$):} Since $P \ge 0$ and $\alpha \in (0, 1)$, we have $\alpha P \ge 0$. Every entry of $\frac{1-\alpha}{n}\IJ_n$ equals $\frac{1-\alpha}{n} > 0$. Therefore, every entry of $G$ satisfies
\[
g_{ij} = \alpha p_{ij} + \frac{1-\alpha}{n} \ge \frac{1-\alpha}{n} > 0
\]
so $G > 0$ is strictly positive (and therefore irreducible and primitive, as $G^1 > 0$)
\vfill
\textbf{2. Row-stochasticity ($G\11_n = \11_n$):} Since $P\11_n = \11_n$ and $\11_n^T \11_n = n$, we compute
\[
G\11_n = \alpha (P\11_n) + \frac{1-\alpha}{n} \11_n (\11_n^T \11_n) = \alpha \11_n + \frac{1-\alpha}{n} n \11_n = (\alpha + 1 - \alpha)\11_n = \11_n
\]
Thus $G$ is row-stochastic, which implies $\rho(G) = 1$
\vfill
\textbf{3. Application of Perron's theorem:} Since $G > 0$ is a positive row-stochastic matrix, Perron's theorem guarantees that $\lambda_1 = \rho(G) = 1$ is simple and strictly dominant, with a unique positive stationary distribution $\bx > \b0$ satisfying $\bx^T G = \bx^T$ and $\11_n^T \bx = 1$, and $\lim_{k\to\infty} (G^T)^k \bx^{(0)} = \bx$ \hfill $\blacksquare$
\end{frame}


%%%%%%%%%%%%%%%%%%%%%%%%%%%%%%%%%
\SSsubsection{Linear system formulation and examples}

\begin{frame}{Reformulating PageRank as a linear system}
Using the normalisation condition $\11_n^T \bx = 1$, we can turn the eigenvalue problem $G^T\bx = \bx$ into a \textbf{non-homogeneous linear system}:
\vfill
\begin{theorem}[Linear system formulation of PageRank]
\label{th:pagerank-linear-system}
Let $P \in \M_n$ be row-stochastic and $\alpha \in (0, 1)$. A vector $\bx \in \IR^n$ with $\11_n^T \bx = 1$ satisfies $\bx^T G = \bx^T$ (or $G^T\bx = \bx$) if and only if it satisfies the linear system
\[
(\II_n - \alpha P^T)\bx = \frac{1 - \alpha}{n}\11_n \qquad \bigl(\text{or } \bx^T(\II_n - \alpha P) = \frac{1-\alpha}{n}\11_n^T\bigr)
\]
Moreover, $\II_n - \alpha P^T$ is invertible with $(\II_n - \alpha P^T)^{-1} = \sum_{k=0}^\infty \alpha^k (P^T)^k$, so the unique PageRank vector is
\[
\bx = \frac{1 - \alpha}{n} (\II_n - \alpha P^T)^{-1} \11_n
\]
\end{theorem}
\end{frame}

\begin{frame}{Proof of Theorem~\ref{th:pagerank-linear-system}}
\textbf{1. Equivalence with $G^T\bx = \bx$:} For any $\bx$ with $\11_n^T \bx = 1$, $\IJ_n^T \bx = \11_n (\11_n^T \bx) = \11_n$, so
\begin{align*}
G^T\bx = \bx
&\iff \alpha P^T \bx + \frac{1-\alpha}{n}\11_n \11_n^T \bx = \bx \\
&\iff \alpha P^T \bx + \frac{1-\alpha}{n}\11_n = \bx \iff (\II_n - \alpha P^T)\bx = \frac{1-\alpha}{n}\11_n
\end{align*}
\textbf{2. Invertibility of $\II_n - \alpha P^T$:} Since $P$ is row-stochastic and $\alpha \in (0, 1)$, $\rho(\alpha P^T) = \alpha \rho(P) = \alpha < 1$, so by Theorem~\ref{th:neumann-series}, $(\II_n - \alpha P^T)^{-1} = \sum_{k=0}^\infty \alpha^k (P^T)^k \ge \b0$
\vfill
\textbf{3. Automatic normalisation ($\11_n^T \bx = 1$):} Left-multiplying by $\11_n^T$ and using $\11_n^T P^T = \11_n^T$ gives
\[
\11_n^T (\II_n - \alpha P^T)\bx = (1 - \alpha)\11_n^T \bx = \frac{1-\alpha}{n}\11_n^T \11_n = 1 - \alpha \implies \11_n^T \bx = 1 \qquad \blacksquare
\]
\end{frame}

\begin{frame}{Unifying thread of Unit 09: the resolvent $(\II - M)^{-1}$}
Notice how all three topics in this unit are governed by the exact same matrix inverse (the \textbf{Neumann series}):
\vfill
\begin{center}
\small
\renewcommand{\arraystretch}{1.6}
\begin{tabular}{l|c|c|c}
\textbf{Topic} & \textbf{Matrix $M$} & \textbf{Condition $\rho(M) < 1$} & \textbf{Resolvent formula} \\ \hline
Absorbing chains & $Q$ & Absorption reachable & $N = (\II_m - Q)^{-1}$ \\
Katz centrality & $\alpha A^T$ & $0 < \alpha < 1/\rho(A)$ & $\bx = \beta(\II_n - \alpha A^T)^{-1}\11_n$ \\
PageRank & $\alpha P^T$ & $0 < \alpha < 1$ & $\bx = \frac{1-\alpha}{n}(\II_n - \alpha P^T)^{-1}\11_n$
\end{tabular}
\end{center}
\vfill
For hand calculations on small graphs ($n = 3$ or $4$), solving the linear system $(\II_n - \alpha P^T)\bx = \frac{1-\alpha}{n}\11_n$ is much faster than forming the dense matrix $G$ and finding its nullspace!
\end{frame}

\begin{frame}{Worked example 1: 3-page web digraph}
Consider a 3-page miniature web where Page $1$ links to Pages $2$ and $3$ ($d_1^+ = 2$), Page $2$ links to Page $3$ ($d_2^+ = 1$) and Page $3$ links to Page $1$ ($d_3^+ = 1$):
\vfill
\begin{columns}[c]
\begin{column}{0.45\textwidth}
\centering
\begin{tikzpicture}[
    ->,
    >=Latex,
    node distance=2.2cm,
    every path/.style={thick},
    pnode/.style={state, fill=myblue!30, minimum size=0.95cm}
]
    \node[pnode] (1) at (0, 1.6) {$1$};
    \node[pnode] (2) at (2.4, 1.6) {$2$};
    \node[pnode] (3) at (1.2, 0) {$3$};
    \draw[->] (1) to node[above] {$1/2$} (2);
    \draw[->] (1) to [bend right=15] node[left] {$1/2$} (3);
    \draw[->] (2) to node[right] {$1$} (3);
    \draw[->] (3) to [bend right=15] node[right] {$1$} (1);
\end{tikzpicture}
\end{column}
\begin{column}{0.55\textwidth}
The row-stochastic transition matrix $P = [p_{ij}]$ (where $p_{ij} = 1/d_i^+$ if $i \to j$) is
\[
P = \begin{pmatrix}
0 & 1/2 & 1/2 \\
0 & 0 & 1 \\
1 & 0 & 0
\end{pmatrix}
\]
With damping factor $\alpha = 1/2$, the Google matrix is
\[
G = \frac{1}{2}P + \frac{1}{6}\IJ_3 = \begin{pmatrix}
1/6 & 5/12 & 5/12 \\
1/6 & 1/6 & 2/3 \\
2/3 & 1/6 & 1/6
\end{pmatrix}
\]
\end{column}
\end{columns}
\end{frame}

\begin{frame}{Worked example 1: setting up $(\II_3 - \alpha P^T)\bx = \frac{1-\alpha}{3}\11_3$}
Let us compute the PageRank vector $\bx$ for $\alpha = 1/2$ by solving $(\II_3 - \frac{1}{2}P^T)\bx = \frac{1}{6}\11_3$:
\vfill
\[
\II_3 - \frac{1}{2}P^T = \begin{pmatrix}
1 & 0 & -1/2 \\
-1/4 & 1 & 0 \\
-1/4 & -1/2 & 1
\end{pmatrix}
\implies
\begin{cases}
x_1 - \frac{1}{2}x_3 = \frac{1}{6} & \text{(Row 1)} \\[4pt]
-\frac{1}{4}x_1 + x_2 = \frac{1}{6} & \text{(Row 2)} \\[4pt]
-\frac{1}{4}x_1 - \frac{1}{2}x_2 + x_3 = \frac{1}{6} & \text{(Row 3)}
\end{cases}
\]
\vfill
From Rows 1 and 2, we can immediately express $x_2$ and $x_3$ in terms of $x_1$:
\[
x_2 = \frac{1}{4}x_1 + \frac{1}{6}, \qquad x_3 = 2x_1 - \frac{1}{3}
\]
\end{frame}

\begin{frame}{Worked example 1: solving for the PageRank vector}
Substituting $x_2 = \frac{1}{4}x_1 + \frac{1}{6}$ and $x_3 = 2x_1 - \frac{1}{3}$ into the normalisation condition $x_1 + x_2 + x_3 = 1$:
\[
x_1 + \left(\frac{1}{4}x_1 + \frac{1}{6}\right) + \left(2x_1 - \frac{1}{3}\right) = 1
\implies \frac{13}{4}x_1 = \frac{7}{6}
\implies x_1 = \frac{14}{39}
\]
\vfill
Back-substituting $x_1 = 14/39$ gives $x_2 = \frac{1}{4}\left(\frac{14}{39}\right) + \frac{1}{6} = \frac{10}{39}$ and $x_3 = 2\left(\frac{14}{39}\right) - \frac{13}{39} = \frac{15}{39}$:
\[
\bx = \begin{pmatrix} 14/39 \\ 10/39 \\ 15/39 \end{pmatrix} \approx \begin{pmatrix} 0.359 \\ 0.256 \\ 0.385 \end{pmatrix} \quad \implies \quad \text{Ranking: Page } 3 \succ \text{Page } 1 \succ \text{Page } 2
\]
\end{frame}

\begin{frame}{Why Page 1 ranks almost as high as Page 3}
Look closely at the ranking in our 3-page example:
\begin{itemize}
\item \textbf{Indegrees:} Page $3$ has indegree $2$ (from $1$ and $2$), while Page $1$ and Page $2$ both have indegree $1$
\item By raw indegree, Page $1$ and Page $2$ would be tied!
\item However, in PageRank, Page $1$ receives the \textbf{undivided vote} of the most important page (Page $3$, with $d_3^+ = 1$), whereas Page $2$ only receives \textbf{half a vote} from Page $1$ ($d_1^+ = 2$)
\item Consequently, Page $1$ ($x_1 = 14/39$) ranks far above Page $2$ ($x_2 = 10/39$) and nearly matches Page $3$ ($x_3 = 15/39$)
\end{itemize}
\vfill
What if we use the standard Google damping factor $\alpha = 0.85$ on the same 3-page network?
\[
(\II_3 - 0.85 P^T)\bx = \frac{0.15}{3}\11_3 \implies \bx \approx \begin{pmatrix} 0.388 \\ 0.215 \\ 0.397 \end{pmatrix}
\]
\end{frame}

\begin{frame}{Worked example 2: 4-page web with a dangling node}
Suppose we add a dangling Page $4$ ($d_4^+ = 0$) to a 4-page web where $1 \to 2$, $1 \to 3$, $2 \to 3$, $3 \to 1$ and $2 \to 4$ (so $d_1^+ = 2$, $d_2^+ = 2$, $d_3^+ = 1$ and $d_4^+ = 0$). Replacing zero row $4$ by $\frac{1}{4}\11_4^T$ gives:
\vfill
\[
P_{\text{raw}} = \begin{pmatrix}
0 & 1/2 & 1/2 & 0 \\
0 & 0 & 1/2 & 1/2 \\
1 & 0 & 0 & 0 \\
0 & 0 & 0 & 0
\end{pmatrix}
\quad \longrightarrow \quad
P = \begin{pmatrix}
0 & 1/2 & 1/2 & 0 \\
0 & 0 & 1/2 & 1/2 \\
1 & 0 & 0 & 0 \\
1/4 & 1/4 & 1/4 & 1/4
\end{pmatrix}
\]
\vfill
Notice why we can solve $(\II_4 - \alpha P^T)\bx = \frac{1-\alpha}{4}\11_4$ even faster using $P_{\text{raw}}$:
\begin{itemize}
\item Write $P = P_{\text{raw}} + \frac{1}{4}\be_4 \11_4^T$, so $P^T\bx = P_{\text{raw}}^T\bx + \frac{x_4}{4}\11_4$
\item Then $(\II_4 - \alpha P_{\text{raw}}^T)\bx = \left(\frac{1-\alpha}{4} + \frac{\alpha x_4}{4}\right)\11_4 = c \11_4$, so we can solve $(\II_4 - \alpha P_{\text{raw}}^T)\by = \11_4$ and normalise $\bx = \by / (\11_4^T \by)$!
\end{itemize}
\end{frame}

\begin{frame}{Computing PageRank at web scale: sparse power iteration}
For the real Web ($n \sim 10^{10}$ pages), how is PageRank computed in practice?
\vfill
\begin{itemize}
\item \textbf{Direct Gaussian elimination} on $(\II_n - \alpha P^T)\bx = \frac{1-\alpha}{n}\11_n$ costs $O(n^3)$ time and $O(n^2)$ memory due to fill-in, which is completely impossible for $n = 10^{10}$
\item \textbf{Forming the dense Google matrix $G$} explicitly also requires $n^2$ entries because $\IJ_n$ is completely dense!
\item \textbf{Matrix-free power iteration:} Starting from $\bx^{(0)} = \frac{1}{n}\11_n$, for any probability vector $\bx^{(k)}$ ($\11_n^T \bx^{(k)} = 1$) we compute
\[
\bx^{(k+1)} = G^T \bx^{(k)} = \alpha P^T \bx^{(k)} + \frac{1 - \alpha}{n}\11_n
\]
\end{itemize}
\vfill
\textbf{Why is this fast?}
\begin{itemize}
\item Each step only requires one sparse matrix-vector product $P^T \bx^{(k)}$ ($O(m)$ time where $m \approx 10n$ is the number of links) plus adding $\frac{1-\alpha}{n}$ to each entry
\item Since $|\lambda_2(G)| \le \alpha = 0.85$, after $k = 50$ steps the error is bounded by $0.85^{50} \approx 3 \times 10^{-4}$
\end{itemize}
\end{frame}

\begin{frame}[fragile]{R implementation: checking our 3-page example}
We verify our hand calculation for $\alpha = 0.5$ in base R using \code{solve}:
\begin{knitrout}
\definecolor{shadecolor}{rgb}{0.969, 0.969, 0.969}\color{fgcolor}\begin{kframe}
\begin{alltt}
\hldef{P3} \hlkwb{<-} \hlkwd{matrix}\hldef{(}\hlkwd{c}\hldef{(}\hlnum{0}\hldef{,} \hlnum{0.5}\hldef{,} \hlnum{0.5}\hldef{,}
               \hlnum{0}\hldef{,} \hlnum{0}\hldef{,}   \hlnum{1}\hldef{,}
               \hlnum{1}\hldef{,} \hlnum{0}\hldef{,}   \hlnum{0}\hldef{),} \hlkwc{nrow} \hldef{=} \hlnum{3}\hldef{,} \hlkwc{byrow} \hldef{=} \hlnum{TRUE}\hldef{)}

\hlcom{# Solve (I - alpha * P^T) x = ((1 - alpha) / n) * 1 for alpha = 0.5}
\hldef{alpha} \hlkwb{<-} \hlnum{0.5}
\hldef{x_half} \hlkwb{<-} \hlkwd{solve}\hldef{(}\hlkwd{diag}\hldef{(}\hlnum{3}\hldef{)} \hlopt{-} \hldef{alpha} \hlopt{*} \hlkwd{t}\hldef{(P3),} \hlkwd{rep}\hldef{((}\hlnum{1} \hlopt{-} \hldef{alpha)} \hlopt{/} \hlnum{3}\hldef{,} \hlnum{3}\hldef{))}
\hldef{MASS}\hlopt{::}\hlkwd{fractions}\hldef{(x_half)}
\end{alltt}
\begin{verbatim}
## [1] 14/39 10/39  5/13
\end{verbatim}
\end{kframe}
\end{knitrout}
\end{frame}

\begin{frame}[fragile]{R implementation: \code{igraph::page\_rank}}
We can also compute PageRank directly from the directed graph using \code{igraph::page\_rank}:
\begin{knitrout}
\definecolor{shadecolor}{rgb}{0.969, 0.969, 0.969}\color{fgcolor}\begin{kframe}
\begin{alltt}
\hlkwd{library}\hldef{(igraph)}
\hldef{edges3} \hlkwb{<-} \hlkwd{matrix}\hldef{(}\hlkwd{c}\hldef{(}\hlnum{1}\hldef{,}\hlnum{2}\hldef{,} \hlnum{1}\hldef{,}\hlnum{3}\hldef{,} \hlnum{2}\hldef{,}\hlnum{3}\hldef{,} \hlnum{3}\hldef{,}\hlnum{1}\hldef{),} \hlkwc{ncol} \hldef{=} \hlnum{2}\hldef{,} \hlkwc{byrow} \hldef{=} \hlnum{TRUE}\hldef{)}
\hldef{g3} \hlkwb{<-} \hlkwd{graph_from_edgelist}\hldef{(edges3,} \hlkwc{directed} \hldef{=} \hlnum{TRUE}\hldef{)}
\hlkwd{page_rank}\hldef{(g3,} \hlkwc{damping} \hldef{=} \hlnum{0.85}\hldef{)}\hlopt{$}\hldef{vector}
\end{alltt}
\begin{verbatim}
## [1] 0.3877897 0.2148106 0.3973997
\end{verbatim}
\end{kframe}
\end{knitrout}
\end{frame}

\begin{frame}[fragile]{R example: 6-page web with a dangling node and spider trap}
Consider a 6-page web graph where Page $4$ is a dangling node ($d_4^+ = 0$) and Pages $\{5, 6\}$ form a spider trap:
\begin{knitrout}
\definecolor{shadecolor}{rgb}{0.969, 0.969, 0.969}\color{fgcolor}\begin{kframe}
\begin{alltt}
\hldef{edges} \hlkwb{<-} \hlkwd{matrix}\hldef{(}\hlkwd{c}\hldef{(}\hlnum{1}\hldef{,}\hlnum{2}\hldef{,} \hlnum{1}\hldef{,}\hlnum{3}\hldef{,} \hlnum{2}\hldef{,}\hlnum{3}\hldef{,} \hlnum{3}\hldef{,}\hlnum{1}\hldef{,} \hlnum{2}\hldef{,}\hlnum{4}\hldef{,} \hlnum{3}\hldef{,}\hlnum{5}\hldef{,} \hlnum{5}\hldef{,}\hlnum{6}\hldef{,} \hlnum{6}\hldef{,}\hlnum{5}\hldef{),}
                \hlkwc{ncol} \hldef{=} \hlnum{2}\hldef{,} \hlkwc{byrow} \hldef{=} \hlnum{TRUE}\hldef{)}
\hldef{g6} \hlkwb{<-} \hlkwd{graph_from_edgelist}\hldef{(edges,} \hlkwc{directed} \hldef{=} \hlnum{TRUE}\hldef{)}
\hlkwd{V}\hldef{(g6)}\hlopt{$}\hldef{name} \hlkwb{<-} \hlkwd{paste0}\hldef{(}\hlsng{"P"}\hldef{,} \hlnum{1}\hlopt{:}\hlnum{6}\hldef{)}
\hldef{pr6} \hlkwb{<-} \hlkwd{page_rank}\hldef{(g6,} \hlkwc{damping} \hldef{=} \hlnum{0.85}\hldef{)}\hlopt{$}\hldef{vector}
\hlkwd{round}\hldef{(pr6,} \hlnum{4}\hldef{)}
\end{alltt}
\begin{verbatim}
##     P1     P2     P3     P4     P5     P6 
## 0.0728 0.0646 0.0920 0.0611 0.3653 0.3442
\end{verbatim}
\begin{alltt}
\hlkwd{par}\hldef{(}\hlkwc{mar} \hldef{=} \hlkwd{c}\hldef{(}\hlnum{0}\hldef{,} \hlnum{0}\hldef{,} \hlnum{0}\hldef{,} \hlnum{0}\hldef{))}
\hlkwd{plot}\hldef{(g6,} \hlkwc{vertex.size} \hldef{=} \hlnum{350} \hlopt{*} \hldef{pr6} \hlopt{+} \hlnum{18}\hldef{,} \hlkwc{vertex.color} \hldef{=} \hlsng{"lightblue"}\hldef{,}
     \hlkwc{vertex.label.cex} \hldef{=} \hlnum{1.1}\hldef{,} \hlkwc{edge.arrow.size} \hldef{=} \hlnum{0.7}\hldef{)}
\end{alltt}
\end{kframe}
\end{knitrout}
\end{frame}

\maxFrameImage{FIGS/U09-U09-pagerank-6page-graph-1.pdf}

\begin{frame}[fragile]{Building the stochastic web matrix in R}
We implement the dangling-node fix and linear system solver on the 6-page example:
\begin{knitrout}
\definecolor{shadecolor}{rgb}{0.969, 0.969, 0.969}\color{fgcolor}\begin{kframe}
\begin{alltt}
\hldef{A6} \hlkwb{<-} \hlkwd{as.matrix}\hldef{(}\hlkwd{as_adjacency_matrix}\hldef{(g6))} \hlcom{# a_\{ij\} = 1 if i -> j}
\hldef{n} \hlkwb{<-} \hlkwd{nrow}\hldef{(A6)}
\hldef{out_deg} \hlkwb{<-} \hlkwd{rowSums}\hldef{(A6)}
\hldef{P6} \hlkwb{<-} \hlkwd{matrix}\hldef{(}\hlnum{0}\hldef{, n, n)}
\hlkwa{for} \hldef{(i} \hlkwa{in} \hlnum{1}\hlopt{:}\hldef{n) \{}
  \hlkwa{if} \hldef{(out_deg[i]} \hlopt{>} \hlnum{0}\hldef{) \{}
    \hldef{P6[i, ]} \hlkwb{<-} \hldef{A6[i, ]} \hlopt{/} \hldef{out_deg[i]}
  \hldef{\}} \hlkwa{else} \hldef{\{}
    \hldef{P6[i, ]} \hlkwb{<-} \hlkwd{rep}\hldef{(}\hlnum{1} \hlopt{/} \hldef{n, n)}
  \hldef{\}}
\hldef{\}}
\hldef{alpha} \hlkwb{<-} \hlnum{0.85}
\hldef{x_exact} \hlkwb{<-} \hlkwd{solve}\hldef{(}\hlkwd{diag}\hldef{(n)} \hlopt{-} \hldef{alpha} \hlopt{*} \hlkwd{t}\hldef{(P6),} \hlkwd{rep}\hldef{((}\hlnum{1} \hlopt{-} \hldef{alpha)} \hlopt{/} \hldef{n, n))}
\end{alltt}
\end{kframe}
\end{knitrout}
\end{frame}

\begin{frame}[fragile]{Power iteration convergence in R}
Now we run 30 steps of sparse power iteration and compare the $\ell_1$ error against $0.85^k$:
\begin{knitrout}
\definecolor{shadecolor}{rgb}{0.969, 0.969, 0.969}\color{fgcolor}\begin{kframe}
\begin{alltt}
\hldef{x_k} \hlkwb{<-} \hlkwd{rep}\hldef{(}\hlnum{1} \hlopt{/} \hldef{n, n)}
\hldef{err} \hlkwb{<-} \hlkwd{numeric}\hldef{(}\hlnum{30}\hldef{)}
\hlkwa{for} \hldef{(k} \hlkwa{in} \hlnum{1}\hlopt{:}\hlnum{30}\hldef{) \{}
  \hldef{x_k} \hlkwb{<-} \hldef{alpha} \hlopt{*} \hlkwd{as.vector}\hldef{(}\hlkwd{t}\hldef{(P6)} \hlopt{%*%} \hldef{x_k)} \hlopt{+} \hldef{(}\hlnum{1} \hlopt{-} \hldef{alpha)} \hlopt{/} \hldef{n}
  \hldef{err[k]} \hlkwb{<-} \hlkwd{sum}\hldef{(}\hlkwd{abs}\hldef{(x_k} \hlopt{-} \hldef{x_exact))}
\hldef{\}}
\hlkwd{plot}\hldef{(}\hlnum{1}\hlopt{:}\hlnum{30}\hldef{, err,} \hlkwc{type} \hldef{=} \hlsng{"b"}\hldef{,} \hlkwc{log} \hldef{=} \hlsng{"y"}\hldef{,} \hlkwc{pch} \hldef{=} \hlnum{19}\hldef{,} \hlkwc{col} \hldef{=} \hlsng{"dodgerblue4"}\hldef{,}
     \hlkwc{xlab} \hldef{=} \hlsng{"Iteration k"}\hldef{,} \hlkwc{ylab} \hldef{=} \hlsng{"L1 error"}\hldef{,}
     \hlkwc{main} \hldef{=} \hlsng{"Geometric convergence of PageRank power iteration"}\hldef{)}
\hlkwd{lines}\hldef{(}\hlnum{1}\hlopt{:}\hlnum{30}\hldef{, err[}\hlnum{1}\hldef{]} \hlopt{*} \hldef{(alpha}\hlopt{^}\hldef{(}\hlnum{0}\hlopt{:}\hlnum{29}\hldef{)),}
      \hlkwc{lty} \hldef{=} \hlnum{2}\hldef{,} \hlkwc{lwd} \hldef{=} \hlnum{2}\hldef{,} \hlkwc{col} \hldef{=} \hlsng{"firebrick"}\hldef{)}
\hlkwd{legend}\hldef{(}\hlsng{"topright"}\hldef{,} \hlkwd{c}\hldef{(}\hlsng{"Actual L1 error"}\hldef{,} \hlsng{"Rate 0.85^k"}\hldef{),}
       \hlkwc{col} \hldef{=} \hlkwd{c}\hldef{(}\hlsng{"dodgerblue4"}\hldef{,} \hlsng{"firebrick"}\hldef{),} \hlkwc{lty} \hldef{=} \hlnum{1}\hlopt{:}\hlnum{2}\hldef{)}
\end{alltt}
\end{kframe}
\end{knitrout}
\end{frame}

\maxFrameImage{FIGS/U09-U09-pagerank-convergence-1.pdf}

\begin{frame}{Summary of Unit 09}
\begin{itemize}
\item \textbf{Absorbing Markov chains:} Reordering states so the $r$ absorbing states precede the $m$ transient states gives the canonical row-stochastic block form
\[
P = \begin{pmatrix} \II_r & \b0_{rm} \\ R & Q \end{pmatrix}, \qquad \rho(Q) < 1
\]
\item \textbf{Fundamental matrix:} $N = (\II_m - Q)^{-1} = \sum_{k=0}^\infty Q^k$ counts expected visits to transient states, $\bt = N\11_m$ gives expected absorption times and $B = N R$ gives absorption probabilities
\item \textbf{Spectral centrality:} Eigenvector centrality ($A^T\bx = \lambda\bx$) requires strong connectivity; Katz centrality ($(\II_n - \alpha A^T)\bx = \beta\11_n$ for $\alpha < 1/\rho(A)$) fixes directed acyclic and source nodes by adding baseline importance
\item \textbf{PageRank:} Normalising votes by outdegree and adding teleportation yields the primitive row-stochastic Google matrix $G = \alpha P + \frac{1-\alpha}{n}\11_n\11_n^T$, whose unique stationary vector solves $(\II_n - \alpha P^T)\bx = \frac{1-\alpha}{n}\11_n$
\end{itemize}
\end{frame}



\end{document}
